\SourcePage{310}
\SGAchapter{XVIII}{整体对偶公式}{La formule de dualité globale}
\begin{center}P. Deligne\end{center}
% 原目录由译本目录重建：引言，曲线的上同调，迹态射，整体对偶定理，参考文献。
\SGAsection{0}{引言}
\textbf{0.1.} 本讲讨论 Poincaré \Term{duality}{对偶}。此\Term{duality}{对偶}依照 Verdier [1] 的模式构造；该文处理\Term{local-scheme}{局部}紧的\Term{hausdorff-space}{分离拓扑空间}。我们先验地证明第 XVII 讲的\Term{functor}{函子} $\mathrm Rf_!$ 有一个\Term{right-adjoint}{右伴随} $\mathrm Rf^!$，并在对\Term{smooth}{光滑}\Term{morphism}{态射}更明确地计算它之前，建立它的许多性质。
这种构造之所以可能，似乎依赖于\Term{functor}{函子} $\mathrm Rf_!$ 的下列性质。

\noindent a) 关于\Term{adjunction}{伴随} $\mathrm Rf^!$ 的存在及其\Term{form}{形式}性质：
(I) \Term{functor}{函子} $\mathrm Rf_!$ 与\Term{base-change}{基变换}交换；
(II) \Term{functor}{函子} $\mathrm R^if_!$ 与\Term{filtered-colimit}{滤过归纳极限}交换，且当 $i$ 充分大时为零；
(III) \Term{functor}{函子} $\mathrm Rf_!$ 可以“在\Term{complex}{复形}层面”计算，即从\Term{complex}{复形}\Term{category}{范畴}之间的一个\Term{functor}{函子}出发计算。

\noindent b) 关于 $f$ \Term{smooth}{光滑}时这个\Term{adjunction}{伴随}的计算：
(IV) 对“小”的 $U$ 计算 $\mathrm Rf_!((\mathbf Z/n\mathbf Z)_U)$。
\SourcePage{311}
\textbf{0.2.} 在独立于第 1、2 节的第 3 节第 1 小节中，我们构造\Term{functor}{函子} $\mathrm Rf^!$；其中使用了本情形特有的构造，而不仅仅使用 (I)、(II)、(III)。这使该小节颇不令人满意；撰稿人承认自己并非总能很好地理解其中究竟发生了什么。
第 2 节讨论 (IV)；一旦有了研究曲线\Term{cohomology}{上同调}的第 1 节，它就是纯\Term{form}{形式}的。
第 1 节还研究\Term{continuous}{连续}系数的情形，因而所建立的结果远多于第 2 节所需。该节深受 Serre [1] 的影响。后文要用到的要点如下：
a) 第 1.1 小节讨论曲线情形的\Term{trace-morphism}{迹态射}，并在 2.9 中推广；
b) 消去定理 1.6.9。愿意接受超越论证的读者，可从 1.6.6（第二个证明）开始阅读第 1.6 小节，而不依赖第 1.2 至 1.5 小节。1.6.9 的关键，是\Term{algebraically-closed-field}{代数闭域}上\Term{smooth}{光滑}曲线的某种\Term{form}{形式}的 Poincaré \Term{duality}{对偶}，以及非循环性引理 XV 2.6（它是\Term{smooth}{光滑}\Term{base-change}{基变换}定理 XVI 1.1 的预备结果）。其结论在 2.17 中推广。
第 3.2 小节收获第 2 节的成果。

\textbf{0.3.} Verdier 的 Poincaré \Term{duality}{对偶}方法允许对\Term{compactifiable}{可紧化}的 $f$ 定义\Term{functor}{函子} $\mathrm Rf^!$。这种一般性的代价，是需要核对一些不常见的相容性，有时即使最简单的相容性，例如 3.2.3，其证明也极其离奇。撰稿人承认，自己证明的相容性比应当证明的少。
有兴趣的读者，可从第 XVII 讲附录出发，毫无困难地将 $\mathrm Rf^!$ 的定义推广到以\Term{coherent}{凝聚}\Term{scheme}{概形}（即\Term{quasi-compact}{拟紧}且\Term{quasi-separated}{拟分离}的\Term{scheme}{概形}）为靶的\Term{finite-type}{有限型}\Term{separated-morphism}{分离态射}。

\textbf{0.4.} 主要结果 3.2.5 用具体的语言表述，因此可以\Term{form}{形式}地返回口头讨论班采用的观点：在那里，\Term{functor}{函子} $\mathrm Rf^!$ 只对\Term{decomposable}{可分解}为一个\Term{smooth}{光滑}\Term{morphism}{态射}与一个\Term{immersion}{浸入}之复合的可\Term{limit}{极限}化（limifiable）\Term{morphism}{态射}定义。口头讨论班证明的梗概可见 Verdier [2]。
\SGAsection{1}{曲线的上同调}
\SGAsection{1.1}{迹态射}
\textbf{1.1.1.} 设 $n\geq1$ 为整数，$X$ 为其上 $n$ \Term{invertible}{可逆}的\Term{scheme}{概形}。$X$ 上的\Term{sheaf}{层} $\mu_n$（IX 3.1）于是是常值\Term{sheaf-of-rings}{环层} $\mathbf Z/n$ 上的秩 $1$ \Term{locally-free}{局部自由}\Term{module-sheaf}{模层}。对 $i\in\mathbf Z$，以 $\mathbf Z/n(i)$ 表示其第 $i$ 次\Term{tensor-power}{张量幂}：
\SourcePage{312}
\begin{equation}\mathbf Z/n(i)=\mathscr Hom(\mathbf Z/n,\mathbf G_m)^{\otimes i}.\tag{1.1.1.1}\end{equation}
若 $n=dn'$，取 $d$ 次幂给出\Term{isomorphism}{同构}
\[\mu_n\otimes_{\mathbf Z/n}\mathbf Z/n'\xrightarrow\sim\mu_{n'},\]
从而得到称为\Term{canonical-isomorphism}{典范同构}的\Term{isomorphism}{同构}
\begin{equation}\mathbf Z/n(i)\otimes_{\mathbf Z/n}\mathbf Z/n'\xrightarrow\sim\mathbf Z/n'(i).\tag{1.1.1.2}\end{equation}
若 $F$ 是满足 $nF=0$ 的\Term{abelian-sheaf}{阿贝尔层}，置
\begin{equation}F(i)=F\otimes_{\mathbf Z/n}\mathbf Z/n(i).\tag{1.1.1.3}\end{equation}
若已有 $n'F=0$，\Term{isomorphism}{同构} (1.1.1.2) 将 $F\otimes_{\mathbf Z/n}\mathbf Z/n(i)$ 与 $F\otimes_{\mathbf Z/n'}\mathbf Z/n'(i)$ 认同，所以 $F(i)$ 只依赖于 $F$，而不依赖于 $n$ 的选取。更一般地，若 $F$ 是与 $X$ 的\Term{residue-characteristic}{剩余特征}互素的\Term{torsion-sheaf}{挠层}（XVII 0.13），并以 ${}_mF$ 表示 $F$ 中乘 $m$ 的核，则定义
\begin{equation}F(i)=\varinjlim_m({}_mF)(i),\tag{1.1.1.4}\end{equation}
其中 $m$ 遍历 $X$ 上\Term{invertible}{可逆}的整数，并在整除关系下趋于无穷。
\Term{functor}{函子} $F\mapsto F(i)$ 是\Term{exact}{正合}的。若 $f:Y\to X$ 是\Term{scheme}{概形}\Term{morphism}{态射}（相应地，是\Term{quasi-compact}{拟紧}\Term{quasi-separated}{拟分离}的\Term{scheme}{概形}\Term{morphism}{态射}），且 $F$ 是 $X$ 上（相应地，$Y$ 上）与 $X$ 的\Term{residue-characteristic}{剩余特征}互素的\Term{torsion-sheaf}{挠层}，则有 $f^*(F(i))=(f^*F)(i)$（相应地，$(f_*F)(i)=f_*(F(i))$）。若 $\mathscr A$ 为 $X$ 上被某个在 $X$ 上\Term{invertible}{可逆}的 $n$ 零化的\Term{sheaf-of-rings}{环层}，则它延拓为 $\mathrm D(X,\mathscr A)$ 到自身的\Term{exact-functor}{正合函子} $K\mapsto K(i)$。这个\Term{functor}{函子}与\Term{inverse-image}{逆像}\Term{functor}{函子}及\Term{direct-image}{直像}\Term{functor}{函子}交换。
\Term{functor}{函子} $F\mapsto F(i)$ 有时称为“\Term{tate-twist}{Tate 扭转}”\Term{functor}{函子}。
\begin{SGAresult}{定义}{1.1.2}
\Term{scheme}{概形} $S$ 上的\Term{flat-curve}{平坦曲线}，是一个\Term{flat}{平坦}、\Term{finitely-presented}{有限表示}、\Term{separated}{分离}且\Term{fiber}{纤维}纯为一维的\Term{morphism}{态射} $f:X\to S$。
\end{SGAresult}
“纯为一维”并不排除空\Term{scheme}{概形}。与 EGA II 7.4.2 不同，这里允许域上的曲线为空，但要求它\Term{separated}{分离}。
当 $S$ 是域的谱时，简称为 $S$ 上的曲线；域上的曲线是\Term{quasi-projective}{拟射影}的。
\SourcePage{313}
\textbf{1.1.3.} 设 $X$ 是\Term{algebraically-closed-field}{代数闭域} $k$ 上的曲线，$k$ 的\Term{characteristic-exponent}{特征指数} $p$ 与整数 $n\geq1$ 互素。先假定 $X$ \Term{reduced}{约化}，取 $k$ 上的\Term{complete}{完备}曲线 $\overline X$，使 $X$ 为其\Term{dense}{稠密}\Term{open-subset}{开子集}。
补集 $Y=\overline X-X$ 是零维的。因此\Term{cohomology}{上同调}\Term{exact-sequence}{正合序列}（XVII 5.1.16.3）\Term{induction-isomorphism}{诱导同构}
\begin{equation}H_c^2(X,\mathbf Z/n(1))\xrightarrow\sim H^2(\overline X,\mathbf Z/n(1)).\tag{1.1.3.1}\end{equation}
还具有由 Kummer 理论给出的\Term{isomorphism}{同构} IX 4.7：
\begin{equation}H^2(\overline X,\mathbf Z/n(1))=\operatorname{Pic}(\overline X)/n=(\mathbf Z/n)^c,\tag{1.1.3.2}\end{equation}
其中 $c$ 是 $\overline X$ 的\Term{irreducible-component}{不可约分支}集，也就是 $X$ 的\Term{irreducible-component}{不可约分支}集。
若不再要求 $X$ \Term{reduced}{约化}，且 $c$ 为其\Term{irreducible-component}{不可约分支}集，则对 $i\in c$，以 $n_i$ 表示 $X$ 在第 $i$ 个\Term{irreducible-component}{不可约分支}的\Term{generic-point}{一般点} $\eta_i$ 处的\Term{multiplicity}{重数}：
\[n_i=\operatorname{lg}(\mathscr O_{X,\eta_i}).\]
已知\Term{restriction}{限制}映射 $H_c^2(X,\mathbf Z/n(1))\to H_c^2(X_{\mathrm{red}},\mathbf Z/n(1))$ 是\Term{isomorphism}{同构}；将 (1.1.3.1)、(1.1.3.2) 用于 $X_{\mathrm{red}}$，便得到 $H_c^2(X,\mathbf Z/n(1))$ 与 $(\mathbf Z/n)^c$ 之间的\Term{canonical-isomorphism}{典范同构}。以 $\operatorname{Tr}_X$，或简称 $\operatorname{Tr}$，表示复合\Term{arrow}{箭头}
\begin{equation}\operatorname{Tr}:H_c^2(X,\mathbf Z/n(1))\xrightarrow\sim(\mathbf Z/n)^c\xrightarrow t\mathbf Z/n,\quad t((a_i)_{i\in c})=\sum n_i a_i.\tag{1.1.3.3}\end{equation}
若 $F$ 是被 $n$ 零化的\Term{abelian-group}{阿贝尔群}，则有\Term{isomorphism}{同构}
\[H_c^2(X,F(1))\xleftrightarrow\sim F\otimes H_c^2(X,\mathbf Z/n(1)),\]
从而由 (1.1.3.3) 得到\Term{trace-morphism}{迹态射}
\begin{equation}\operatorname{Tr}:H_c^2(X,F(1))\longrightarrow F.\tag{1.1.3.4}\end{equation}
若 $n=n'd$，则图
\begin{equation}\begin{tikzcd}
0\ar[r]&\mu_n\ar[r]\ar[d,"{x^d}"]&\mathbf G_m\ar[r,"{x^n}"]\ar[d,"{x^d}"]&\mathbf G_m\ar[r]\ar[d,equal]&0\\
0\ar[r]&\mu_{n'}\ar[r]&\mathbf G_m\ar[r,"{x^{n'}}"]&\mathbf G_m\ar[r]&0
\end{tikzcd}\tag{1.1.3.5}\end{equation}
\SourcePage{314}
交换；根据\Term{isomorphism}{同构} (1.1.3.2) 在 IX 4.7 中的定义，以下图也交换：
\begin{equation}\begin{tikzcd}H_c^2(X,\mathbf Z/n(1))\ar[r,"\sim"]\ar[d]&(\mathbf Z/n)^c\ar[d]\\H_c^2(X,\mathbf Z/n'(1))\ar[r,"\sim"]&(\mathbf Z/n')^c\end{tikzcd}\tag{1.1.3.6}\end{equation}
所以\Term{arrow}{箭头} (1.1.3.4) 只依赖于 $F$，而不依赖于满足 $nF=0$ 且与 $p$ 互素的整数 $n$ 的选取。取\Term{limit}{极限}，便对一切与 $p$ 互素的挠\Term{abelian-group}{阿贝尔群} $F$ 定义了它。
以下引理直接来自定义。
\begin{SGAresult}{引理}{1.1.4}
在上述假设下，若对 $i\in c$，$U_i$ 是包含于 $X$ 的第 $i$ 个\Term{irreducible-component}{不可约分支}中的\Term{nonempty}{非空}\Term{open-subset}{开子集}，则下图交换：
\[\begin{tikzcd}\displaystyle\bigoplus_{i\in c}H_c^2(U_i,F(1))\ar[rr,"\sim"]\ar[dr,"{\bigoplus_{i\in c}\operatorname{Tr}_{U_i}}"']&&H_c^2(X,F(1))\ar[dl,"\operatorname{Tr}_X"]\\&F&\end{tikzcd}\]
\end{SGAresult}
\begin{SGAresult}{引理}{1.1.5}
设 $X,Y$ 为\Term{algebraically-closed-field}{代数闭域} $k$ 上的两条曲线，$u:X\to Y$ 为\Term{quasi-finite}{拟有限}\Term{flat}{平坦}\Term{morphism}{态射}，$n$ 为与 $k$ 的\Term{characteristic-exponent}{特征指数} $p$ 互素的整数。则下图交换，其中 $\operatorname{Tr}_u$ 为\Term{arrow}{箭头} XVII 6.2.3：
\[\begin{tikzcd}H_c^2(X,\mathbf Z/n(1))\ar[rr,"\operatorname{Tr}_u"]\ar[dr,"\operatorname{Tr}_X"']&&H_c^2(Y,\mathbf Z/n(1))\ar[dl,"\operatorname{Tr}_Y"]\\&\mathbf Z/n&\end{tikzcd}\]
\end{SGAresult}
\begin{proof}
由引理 1.1.4，可归结于 $X,Y$ \Term{irreducible}{不可约}且 $X_{\mathrm{red}},Y_{\mathrm{red}}$ \Term{smooth}{光滑}的情形。若 $e$（相应地 $f$）是 $X$（相应地 $Y$）在\Term{generic-point}{一般点}处的\Term{multiplicity}{重数}，
\SourcePage{315}
且 $e=df$，则以下图中标记 $+$ 的三角形交换：
\[\begin{tikzcd}[column sep=large]H_c^2(X,\mathbf Z/n(1))\ar[rr,"\sim"]\ar[dd,"\operatorname{Tr}_u"']\ar[dr,"\operatorname{Tr}_X"]&&H_c^2(X_{\mathrm{red}},\mathbf Z/n(1))\ar[dd,"{d\operatorname{Tr}_{u_{\mathrm{red}}}}"]\ar[dl,"{e\operatorname{Tr}_{X_{\mathrm{red}}}}"']\\&\mathbf Z/n&\\H_c^2(Y,\mathbf Z/n(1))\ar[rr,"\sim"]\ar[ur,"\operatorname{Tr}_Y"]&&H_c^2(Y_{\mathrm{red}},\mathbf Z/n(1))\ar[ul,"{f\operatorname{Tr}_{Y_{\mathrm{red}}}}"']\end{tikzcd}\]
其外框也交换，这可从以下\Term{commutative-diagram}{交换图}推出：
\[\begin{tikzcd}u_!\mathbf Z/n_X(1)\ar[r,"\sim"]\ar[d,"\operatorname{Tr}_u"']&u_{\mathrm{red}!}\mathbf Z/n_{X_{\mathrm{red}}}(1)\ar[d,"{d\operatorname{Tr}_{u_{\mathrm{red}}}}"]\\\mathbf Z/n_Y(1)\ar[r,"\sim"]&\mathbf Z/n_{Y_{\mathrm{red}}}(1).\end{tikzcd}\]
只需在 $Y$ 的\Term{generic-point}{一般点}验证后一相容性，而这是显然的。
设 $\overline X,\overline Y$ 是分别以 $X_{\mathrm{red}},Y_{\mathrm{red}}$ 为\Term{dense}{稠密}\Term{open-subset}{开子集}的非\Term{smooth}{光滑}\Term{complete}{完备}曲线（EGA II 7.4.11）。\Term{morphism}{态射} $u_{\mathrm{red}}$ 延拓为\Term{flat}{平坦}\Term{morphism}{态射} $\overline u:\overline X\to\overline Y$（EGA II 7.4.9）；由 1.1.4，只需对 $\overline u$ 验证 1.1.5。
由 XVII 6.3.18.2，以下 $Y$ 上的\Term{sheaf}{层}图交换，其中 $N$ 表示\Term{norm}{范数}：
\[\begin{tikzcd}0\ar[r]&\overline u_*\mu_n\ar[r]\ar[d,"\operatorname{Tr}_u"]&\overline u_*\mathbf G_m\ar[r,"{x^n}"]\ar[d,"N"]&\overline u_*\mathbf G_m\ar[r]\ar[d,"N"]&0\\0\ar[r]&\mu_n\ar[r]&\mathbf G_m\ar[r,"{x^n}"]&\mathbf G_m\ar[r]&0.\end{tikzcd}\]
因此以下图也交换：
\[\begin{tikzcd}\operatorname{Pic}(X)/n\ar[r,"\sim"]\ar[d,"N"']&H_c^2(X,\mathbf Z/n(1))\ar[d,"\operatorname{Tr}_u"]\\\operatorname{Pic}(Y)/n\ar[r,"\sim"]&H_c^2(Y,\mathbf Z/n(1)).\end{tikzcd}\]
剩下只须指出：若 $\mathscr L$ 是 $X$ 上次数为 $1$ 的\Term{invertible-sheaf}{可逆层}，则其\Term{norm}{范数}
\SourcePage{316}
的次数仍为 $1$（这由相容性 EGA IV 21.10.7 得到，同时利用 EGA IV 21.5.5 对\Term{divisor}{除子}\Term{norm}{范数}的定义）。
\end{proof}
\begin{SGAresult}{命题}{1.1.6}
对于每条\Term{compactifiable}{可紧化}的\Term{flat-curve}{平坦曲线} $f:X\to S$ 和每个 $S$ 上与 $S$ 的\Term{residue-characteristic}{剩余特征}互素的\Term{torsion-sheaf}{挠层} $F$，存在唯一的方法定义\Term{trace-morphism}{迹态射}
\[\operatorname{Tr}_f:\mathrm R^2f_!(f^*F(1))\longrightarrow F,\]
使之关于 $F$ \Term{functorial}{函子性}，构造与任意\Term{base-change}{基变换}相容（参见 XVII 6.2.3，(VAR 2)），且当 $S$ 是\Term{algebraically-closed-field}{代数闭域}的谱时与\Term{trace-morphism}{迹态射} (1.1.3.4) 一致。
\end{SGAresult}
\begin{proof}
\Term{morphism}{态射} $\operatorname{Tr}_f$ 已在每个\Term{fiber}{纤维}上给定，故唯一性是清楚的。
设 $p:\mathbf P^1_S\to S$ 为 $S$ 上的\Term{projective}{射影}直线，$n$ 是 $S$ 上\Term{invertible}{可逆}的整数。有 $\operatorname{Pic}_S(\mathbf P^1_S)=\mathbf Z_S$，因而由 Kummer 理论（IX 3.2）得到\Term{morphism}{态射} $\mathbf Z/n\to\mathrm R^2p_*\mathbf Z/n(1)$，逐\Term{fiber}{纤维}验证（XII 5.2）可知这是\Term{isomorphism}{同构}。其逆 $\operatorname{Tr}_p$ 在每个\Term{fiber}{纤维}上诱导\Term{trace-morphism}{迹态射} (1.1.3.4)。
设 $f:X\to S$ 为 $S$ 上的\Term{flat-curve}{平坦曲线}，并存在\Term{quasi-finite}{拟有限}\Term{flat}{平坦}的 $S$-\Term{morphism}{态射} $u:X\to\mathbf P^1_S$。由 1.1.5，\Term{composition-morphism}{复合态射} $\operatorname{Tr}_p\circ\operatorname{Tr}_u$（XVII 6.2.3）在 $S$ 的任一\Term{geometric-point}{几何点}处的\Term{fiber}{纤维}为\Term{trace-morphism}{迹态射} (1.1.3.4)。特别地，它不依赖 $u$ 的选择。假设 $X$ 是一族开子\Term{scheme}{概形} $U_i$ 的并，且每个 $U_i$ 都有到 $\mathbf P^1_S$ 的\Term{quasi-finite}{拟有限}\Term{flat}{平坦}\Term{morphism}{态射}。若 $\alpha_i$（相应地 $\alpha_{ij}$）为 $U_i$（相应地 $U_{ij}=U_i\cap U_j$）到 $X$ 的包含，则序列
\[\bigoplus\mathrm R^2(f\alpha_{ij})_!\mathbf Z/n(1)\rightrightarrows\bigoplus\mathrm R^2(f\alpha_i)_!\mathbf Z/n(1)\longrightarrow\mathrm R^2f_!\mathbf Z/n(1)\longrightarrow0\]
由 XVII 6.2.8 及 $\mathrm R^2f_!$ 的\Term{right-exact}{右正合}性（因为 $\mathrm R^3f_!=0$）而\Term{exact}{正合}。利用 1.1.5 逐\Term{fiber}{纤维}验证，\Term{trace-morphism}{迹态射}之和 $\sum_i\mathrm R^2(f\alpha_i)_!\mathbf Z/n(1)\to\mathbf Z/n$ 通过 $\mathrm R^2f_!\mathbf Z/n(1)$ 的一个\Term{trace-morphism}{迹态射} $\operatorname{Tr}_f$ \Term{resolution}{分解}，该\Term{morphism}{态射}在每个\Term{fiber}{纤维}上诱导 (1.1.3.4)。
对一般的 $X$，设 $j:X'\to X$ 为使 $f|_{X'}$ 成为 Cohen--Macaulay \Term{morphism}{态射}的最大\Term{open-subset}{开子集}（参见 EGA IV 12.1.1 (vi)）。$X'$ 的每点 $x$ 在 $X'$ 中都有\Term{neighborhood}{邻域} $U_x$，以及\Term{quasi-finite}{拟有限}的 $S$-\Term{morphism}{态射} $u_x:U_x\to\mathbf P^1_S$。此\Term{morphism}{态射}自动\Term{flat}{平坦}（EGA IV 11.3.10、15.4.2），故 $X'$ 满足前述假设。此外 $X'$ 在 $X$ 的每个\Term{fiber}{纤维}中\Term{dense}{稠密}，所以逐\Term{fiber}{纤维}可见 $\mathrm R^2(fj)_!\mathbf Z/n(1)\to\mathrm R^2f_!\mathbf Z/n(1)$ 为\Term{isomorphism}{同构}；复合\Term{arrow}{箭头}
\[\operatorname{Tr}_f:\mathrm R^2f_!\mathbf Z/n(1)\xleftarrow\sim\mathrm R^2(fj)_!\mathbf Z/n(1)\xrightarrow{\operatorname{Tr}_{fj}}\mathbf Z/n\]
在每个\Term{fiber}{纤维}上诱导 (1.1.3.4)。
若 $nF=0$，由 XVII 5.2.6 有
\begin{equation}\mathrm R^2f_!\mathbf Z/n(1)\otimes F\xrightarrow\sim\mathrm R^2f_!(f^*F(1)),\tag{1.1.6.1}\end{equation}
\SourcePage{317}
并将 $\operatorname{Tr}_f$ 定义为复合\Term{arrow}{箭头}
\[\operatorname{Tr}_f:\mathrm R^2f_!(f^*F(1))\xleftarrow\sim\mathrm R^2f_!\mathbf Z/n(1)\otimes F\longrightarrow F.\]
在一般情形，以 $S$ 上\Term{invertible}{可逆}的 $n$ 所对应的\Term{sheaf}{层} ${}_nF$ 的\Term{trace-morphism}{迹态射}之\Term{inductive-limit}{归纳极限}定义 $\operatorname{Tr}_f$。这些\Term{arrow}{箭头}在每个\Term{fiber}{纤维}上诱导 (1.1.3.4)，故与任意\Term{base-change}{基变换}交换。
\end{proof}
利用 1.1.4、1.1.5 和 (1.1.6.1)，逐\Term{fiber}{纤维}验证可得下列结果。
\begin{SGAresult}{引理}{1.1.7}
设 $f:Y\to S$ 是 $S$ 上\Term{compactifiable}{可紧化}的\Term{flat-curve}{平坦曲线}，$u:X\to Y$ 是\Term{finitely-presented}{有限表示}的\Term{quasi-finite}{拟有限}\Term{flat}{平坦}\Term{morphism}{态射}，$F$ 是 $S$ 上与 $S$ 的\Term{residue-characteristic}{剩余特征}互素的\Term{torsion-sheaf}{挠层}。则下图交换：
\[\begin{tikzcd}\mathrm R^2(fu)_!(fu)^*F(1)\ar[r,equal]\ar[d,"\operatorname{Tr}_{fu}"']&\mathrm R^2f_!(u_!u^*)f^*F(1)\ar[r,"\operatorname{Tr}_u"]&\mathrm R^2f_!f^*F(1)\ar[d,"\operatorname{Tr}_f"]\\F\ar[rr,equal]&&F.\end{tikzcd}\]
\end{SGAresult}
\begin{SGAresult}{引理}{1.1.8}
设 $u:Y\to S$ 是\Term{finitely-presented}{有限表示}的\Term{quasi-finite}{拟有限}\Term{flat}{平坦}\Term{morphism}{态射}，$f:X\to Y$ 是 $S$-\Term{compactifiable}{可紧化}的\Term{flat-curve}{平坦曲线}（XVII 3.2.1），$F$ 如 1.1.7。则下图交换：
\[\begin{tikzcd}\mathrm R^2(uf)_!(uf)^*F(1)\ar[r,equal]\ar[d,"\operatorname{Tr}_{uf}"']&u_!\mathrm R^2f_!f^*u^*F(1)\ar[r,"\operatorname{Tr}_f"]&u_!u^*F\ar[d,"\operatorname{Tr}_u"]\\F\ar[rr,equal]&&F.\end{tikzcd}\]
\end{SGAresult}
同样，由 (1.1.3.2) 验证可得：
\begin{SGAresult}{引理}{1.1.9}
若 $f:X\to S$ 是\Term{compactifiable}{可紧化}的\Term{smooth}{光滑}曲线，且 $f$ 的\Term{geometric-fiber}{几何纤维}\Term{irreducible}{不可约}，则\Term{trace-morphism}{迹态射} (1.1.6) 是\Term{isomorphism}{同构}。
\end{SGAresult}
\SGAsection{1.2}{射影空间的 $1$-非循环性}
本小节原则上始终使用 fppf \Term{topology}{拓扑}（XVII 0.10）。
\textbf{1.2.1.} 设 $\mathscr E$ 为\Term{scheme}{概形} $S$ 上处处秩 $\geq2$ 的\Term{locally-free}{局部自由}模，$p:\mathbf P(\mathscr E)\to S$ 为相应的\Term{projective-bundle}{射影丛}，$G$ 为 $S$ 的大 fppf 景上的\Term{abelian-sheaf}{阿贝尔层}（XVII 0.10）。每个二元组 $(K_0,\varphi)$，其中 $K_0$ 为 $G$ 下的\Term{torsor}{挠子}，$\varphi:\mathbf G_{mS}\to G$ 为\Term{homomorphism}{同态}，都定义 $\mathbf P(\mathscr E)$ 上的一个 $p^*G$-\Term{torsor}{挠子} $K(K_0,\varphi)$，即 $K_0$ 与 $\mathbf G_m$-\Term{torsor}{挠子} $\mathscr O(1)$ 在 $\varphi$ 下的像之和：
\begin{equation}K(K_0,\varphi)=K_0+\varphi\mathscr O(1).\tag{1.2.1.1}\end{equation}
\SourcePage{318}
令
\[\operatorname{Hom}((K_0,\varphi),(K'_0,\varphi'))=\begin{cases}\varnothing&\varphi\ne\varphi',\\\operatorname{Hom}(K_0,K'_0)&\varphi=\varphi',\end{cases}\]
便使这些二元组成为一个\Term{category}{范畴}。
于是\Term{torsor}{挠子} $K(K_0,\varphi)$ \Term{functorial}{函子性}地依赖于 $(K_0,\varphi)$。这一构造与\Term{localization}{局部化}相容，从而定义了从 $S$ 上二元组 $(K_0,\varphi)$ 构成的\Term{stack}{叠}到 $S$ 上 $\mathbf P(\mathscr E)$ 上的 $p^*G$-\Term{torsor}{挠子}所构成的\Term{stack}{叠}的\Term{functor}{函子}（后者在 $U/S$ 上的\Term{section}{截面}就是 $U\times_S\mathbf P(\mathscr E)$ 上的 $p^*G$-\Term{torsor}{挠子}）。这些\Term{stack}{叠}都是\Term{groupoid}{群胚}\Term{stack}{叠}（即任一 $U/S$ 上方的每个\Term{morphism}{态射}都是\Term{isomorphism}{同构}）。
\begin{SGAresult}{定理}{1.2.2}
在 1.2.1 的假设下，若 $G$ 为 $S$ 上\Term{flat}{平坦}、\Term{finitely-presented}{有限表示}的交换\Term{group-scheme}{群概形}，则 1.2.1 的\Term{functor}{函子} $K$ 为\Term{equivalence}{等价}。
\end{SGAresult}
\textbf{1.2.2.1.} 此陈述\Term{equivalence}{等价}于以下两式同时成立：
\[G\xrightarrow\sim p_*p^*G,\qquad\mathscr Hom(\mathbf G_m,G)\xrightarrow\sim\mathrm R^1p_*p^*G.\]
事实上，一个\Term{groupoid}{群胚}\Term{stack}{叠}\Term{morphism}{态射}为\Term{equivalence}{等价}，当且仅当它在由\Term{object}{对象}\Term{isomorphism-class}{同构类}的\Term{presheaf}{预层}生成的\Term{sheaf}{层}上\Term{induction-isomorphism}{诱导同构}（第二式），并在任一\Term{local-scheme}{局部}\Term{object}{对象}的\Term{automorphism}{自同构}\Term{sheaf}{层}上\Term{induction-isomorphism}{诱导同构}（第一式）。
我们将给出\Term{fully-faithful}{完全忠实}性的两个证明（1.2.4、1.2.5）。
\begin{SGAresult}{引理}{1.2.3}
设 $f:X\to S$ 为\Term{proper}{固有}\Term{flat}{平坦}\Term{morphism}{态射}。假设\Term{universal}{普遍}地有 $f_*\mathscr O_X=\mathscr O_S$ 及 $\mathrm R^1f_*\mathscr O_X=0$。则每个从 $X$ 到\Term{finite-type}{有限型} $S$-\Term{group-scheme}{群概形} $G$ 的 $S$-\Term{morphism}{态射}，都通过 $f$ 和 $G$ 的一个\Term{section}{截面}\Term{resolution}{分解}；更确切地，$G\xrightarrow\sim f_*f^*G$。
\end{SGAresult}
\begin{proof}
只须证明第一个断言。根据\Term{rigidity}{刚性}引理（Mumford [1, 6.1, p. 115]；注意 Mumford 并未使用对 $S$ 所作的\Term{noetherian}{诺特}假设），只须处理 $S$ 是\Term{algebraically-closed-field}{代数闭域}的谱的情形。先假设 $G$ 为\Term{abelian-scheme}{阿贝尔概形}，以 $G^*$ 为其\Term{duality}{对偶}\Term{abelian-scheme}{阿贝尔概形}。因为 $G=G^{**}$，给出 $\varphi:X\to G$ \Term{equivalence}{等价}于给出 $X\times G^*$ 上的\Term{invertible-sheaf}{可逆层} $\mathscr L$，使它沿 $X\times\{e\}$ \Term{trivialization}{平凡化}，并在\Term{projection}{投影} $\operatorname{pr}_1:X\times G^*\to X$ 的\Term{fiber}{纤维}上\Term{algebra}{代数}\Term{equivalence}{等价}于零。若 $x:S\to X$ 为 $X$ 的\Term{rational-point}{有理点}，则 $\mathscr L\simeq\operatorname{pr}_2^*(x_{G^*})^*\mathscr L\otimes\mathscr M$，其中 $\mathscr M$ 沿 $x\times G^*$ 及 $X\times\{e\}$ \Term{trivialization}{平凡化}。给出 $\mathscr M$ \Term{equivalence}{等价}于给出一个将 $e$ 映到 $e$ 的\Term{scheme}{概形}\Term{morphism}{态射} $G^*\to\operatorname{Pic}_{X/S}$。因 $H^1(X,\mathscr O_X)=0$，有 $\operatorname{Pic}^0_{X/S}=0$，$\mathscr M\simeq\mathscr O$，而 $\mathscr L\simeq\operatorname{pr}_2^*\mathscr L'$，其中 $\mathscr L'$ 在 $G^*$ 上\Term{equivalence}{等价}于零，因而定义了满足 $\varphi=g\circ f$ 的 $g\in G(S)$。
\SourcePage{319}
一般情形下，$G$ 是\Term{abelian-scheme}{阿贝尔概形} $A$ 被\Term{affine}{仿射}群 $G_0$ 的\Term{group-extension}{扩张}。上述结果使我们归结于 $A=0$，此时有
\[\operatorname{Hom}_S(\operatorname{Spec}H^0(X,\mathscr O_X),G)\xrightarrow\sim\operatorname{Hom}_S(X,G).\]
\end{proof}
\begin{SGAresult}{引理}{1.2.4}
在 1.2.1 的假设下，若 $G$ 是 $S$ 上\Term{finitely-presented}{有限表示}的\Term{group-scheme}{群概形}，或由 $S$ 上的\Term{quasi-coherent}{拟凝聚}\Term{sheaf}{层}定义，则\Term{functor}{函子} (1.2.1.1) \Term{fully-faithful}{完全忠实}。
\end{SGAresult}
\begin{proof}
A. 证明若 $K(K_0,\varphi)$ \Term{isomorphism}{同构}于 $K(K'_0,\varphi')$，则 $\varphi=\varphi'$。当 $G$ \Term{quasi-coherent}{拟凝聚}时，$\varphi=\varphi'=0$，断言无须证明。当 $G$ \Term{finitely-presented}{有限表示}时，由 SGA 3 IX 5.1，只须对 $S$ 为\Term{algebraically-closed-field}{代数闭域} $k$ 的谱证明此断言。设 $T$ 为 $G$ 的最大子\Term{torus}{环面}，$H=G/T$。\Term{cohomology}{上同调}\Term{exact-sequence}{正合序列}给出
\[0\to H^0(\mathbf P(\mathscr E),T)\to H^0(\mathbf P(\mathscr E),G)\xrightarrow\alpha H^0(\mathbf P(\mathscr E),H)\xrightarrow\delta H^1(\mathbf P(\mathscr E),T)\xrightarrow\beta H^1(\mathbf P(\mathscr E),G).\]
由 1.2.3，$H^0(\mathbf P(\mathscr E),G)\simeq G(k)$ 且 $H^0(\mathbf P(\mathscr E),H)\simeq H(k)$，所以 $\alpha$ \Term{surjectivity}{满射}，$\delta$ 为零。若 $Y(T)=\operatorname{Hom}(\mathbf G_m,T)$，则 $T\simeq\mathbf G_m\otimes Y(T)$，$H^1(\mathbf P(\mathscr E),T)\simeq Y(T)$。因此 $\beta$ 的\Term{injectivity}{单射}性表达了由 $\mathscr O(1)$ 定义的\Term{canonical}{典范}\Term{arrow}{箭头}
\[\operatorname{Hom}(\mathbf G_m,G)\simeq\operatorname{Hom}(\mathbf G_m,T)\longrightarrow H^1(\mathbf P(\mathscr E),G)\]
的\Term{injectivity}{单射}性，从而证明了断言。
B. \Term{functor}{函子} (1.2.1.1) 的\Term{fully-faithful}{完全忠实}性是 $S$ 上的\Term{local-scheme}{局部}问题。因此可假设 $\mathbf P(\mathscr E)/S$ 有\Term{section}{截面} $s$。若 $K(K_0,\varphi)\simeq K(K'_0,\varphi')$，则由 A 有 $\varphi=\varphi'$，以及
$K_0\simeq s^*K(K_0,\varphi)\simeq s^*K(K'_0,\varphi')\simeq K'_0$。
利用公式
\begin{equation}\operatorname{Hom}(K(K_0,\varphi),K(K'_0,\varphi'))\simeq\operatorname{Hom}(K(K_0-K'_0,\varphi-\varphi'),G_{\mathbf P(\mathscr E)}),\tag{1.2.4.1}\end{equation}
只剩下证明
\[\operatorname{Hom}_{\mathrm{tors}}(G,G)\xrightarrow\sim\operatorname{Hom}_{\mathrm{tors}}(p^*G,p^*G),\]
这由 1.2.3 得到。
\end{proof}
\textbf{1.2.5.} 下面是\Term{fully-faithful}{完全忠实}性的另一个、更具 Künneth 风格的证明。问题在 $S$ 上是\Term{local-scheme}{局部}的；由 (1.2.4.1)，只须验证
\[\operatorname{Hom}_{\mathbf P(\mathscr E)}(G,K(G,\varphi))\begin{cases}=\varnothing&\varphi\ne0,\\\xleftarrow\sim G&\varphi\ne0.\end{cases}\]
这正是以下引理所表达的。
\begin{SGAresult}{引理}{1.2.6}
设 $\mathscr E$ 为 $S$ 上处处秩 $\geq2$ 的\Term{locally-free}{局部自由}模，$G$ 为 $S$ 上\Term{finitely-presented}{有限表示}的\Term{group-scheme}{群概形}，$\varphi:\mathbf G_m\to G$ 为\Term{group-homomorphism}{群同态}，$\mathbf V(\mathscr E)$ 为 $\mathscr E$ 定义的\Term{vector-bundle}{向量丛}，$\mathbf V(\mathscr E)^*$ 为相应的去零\Term{vector-bundle}{向量丛}，$\Psi:\mathbf V(\mathscr E)^*\to G$ 为满足以下 $\mathbf G_m\times\mathbf V(\mathscr E)^*\to G$ 的 $S$-\Term{morphism}{态射}恒等式的 $S$-\Term{morphism}{态射}：
\[\Psi(\lambda v)=\varphi(\lambda)\Psi(v).\]
\SourcePage{320}
则 $\varphi=e$，且 $\Psi$ 通过 $\mathbf P(\mathscr E)$ 到 $S$ 的\Term{projection}{投影} $p$ \Term{resolution}{分解}。
\end{SGAresult}
\begin{proof}
只须证明第二个断言。立刻可归结于 $S$ \Term{noetherian}{诺特}的情形，然后由 SGA 3 IX 5.1 及应用于 $\mathbf P(\mathscr E)$ 的\Term{rigidity}{刚性}引理（Mumford [1, 6.1, p. 115]），归结于 $S$ 为域的谱。设 $n$ 为整数，$W_n,G_n$ 为 $\mathbf V(\mathscr E),G$ 中\Term{zero-section}{零截面}的第 $n$ 阶\Term{infinitesimal-neighborhood}{无穷小邻域}，$g_n$ 为 $G_n$ 的\Term{affine}{仿射}\Term{algebra}{代数}，$P_n(\Psi):\mathbf V(\mathscr E)^*\to\mathscr Hom_S(W_n,G)$ 为 $\Psi$ 的第 $n$ 阶主部。$\mathscr Hom_S(W_n,G)$ 的每个\Term{section}{截面}，通过按\Term{zero-section}{零截面}像的逆作右\Term{translation}{平移}，定义一个将 $0$ 映到 $e$ 的 $W_n\to G_n$ 的映射；这定义了
\[\Psi_1=P_n(\Psi)\Psi^{-1}:\mathbf V(\mathscr E)^*\longrightarrow\mathscr Hom_S(W_n,G_n).\]
假设意味着 $\Psi_1$ 与\Term{homothety}{位似}交换；\Term{homothety}{位似}在右侧的\Term{action}{作用}，由其在 $W_n\subset\mathbf V(\mathscr E)$ 上的\Term{action}{作用}搬运结构得到。$W_n$ 的\Term{affine}{仿射}\Term{algebra}{代数}为 $\sum_{i=0}^n\operatorname{Sym}^i(\mathscr E)$。由\Term{homogeneous}{齐次}性，$\Psi_1$ 由 $\mathbf P(\mathscr E)$ 上 $\mathscr O_{\mathbf P(\mathscr E)}$-模
\[\sum_{i=0}^np^*\mathscr Hom(g_n,\operatorname{Sym}^i(\mathscr E))(-i)\]
的一个\Term{section}{截面}定义。这个模中指标 $i>0$ 的\Term{set-component}{分量}只有\Term{zero-section}{零截面}，所以 $\Psi_1=0$。
对每个 $n$ 都如此，故 $\Psi$ 通过 $G$ 的一个\Term{section}{截面}\Term{resolution}{分解}。
\end{proof}
这第二个证明还可推广，用来证明下面的 1.2.4 的解析变式。当 $G$ \Term{affine}{仿射}时，1.2.2 的证明可以独立于该变式进行。
\textbf{1.2.7.} 设 $L$ 为\Term{complete-discrete-valuation-ring}{完备离散赋值环} $\mathscr O$ 的\Term{field-of-fractions}{分式域}，$U_m$ 为 $\mathbf G_m$ 的\Term{formal-completion}{形式完备化}之“通有\Term{fiber}{纤维}”\Term{rigidity}{刚性}解析群，$\mathscr O(1)^\wedge$ 为 $\mathbf P^r_L$ 上作为 $\mathscr O(1)$ 的\Term{formal-completion}{形式完备化}的 $U_m$-\Term{torsor}{挠子}。若 $G$ 为 $L$ 上的\Term{rigidity}{刚性}解析群，则\Term{functor}{函子}
\[(K_0,\varphi)\longmapsto p^*K_0+\varphi\mathscr O(1)^\wedge\]
（参见 1.2.1）从由 $\operatorname{Spec}(L)$ 上的 $G$-\Term{torsor}{挠子}及 $\varphi\in\operatorname{Hom}(U,G)$ 组成的二元组 $(K_0,\varphi)$ 的\Term{category}{范畴}，到 $\mathbf P^r_L$ 上的 $p^*G$-\Term{torsor}{挠子}\Term{category}{范畴}，是\Term{fully-faithful}{完全忠实}的。
\begin{SGAresult}{引理}{1.2.8}
设 $k$ 为\Term{characteristic-exponent}{特征指数} $p$ 的域，$f:X\to\operatorname{Spec}(k)$ 为 $k$ 上满足 $H^0(X,\mathscr O_X)=k$ 及 $\operatorname{Pic}^\tau_X=0$ 的\Term{proper}{固有}\Term{scheme}{概形}。对每个 $k$ 上不含\Term{torus}{环面}的\Term{finite-type}{有限型}\Term{affine}{仿射}交换\Term{group-scheme}{群概形} $G$，有 $\mathrm R^1f_*(f^*G)=0$。
\end{SGAresult}
\begin{proof}
只须证明在任意\Term{base-change}{基变换} $c:S\to\operatorname{Spec}(k)$ 后，$f^*$ 从 $S$ 上的 $G$-\Term{torsor}{挠子}\Term{category}{范畴}到 $X_S$ 上的 $G$-\Term{torsor}{挠子}\Term{category}{范畴}为\Term{equivalence}{等价}。问题在 $\operatorname{Spec}(k)$ 上是\Term{local-scheme}{局部}的，因此可归结于 $X$ 有一个\Term{rational-point}{有理点}，$G$ 是群 $\mathbf G_a,\alpha_p,\mu_p,\mathbf Z/p,\mathbf Z/\ell_i$ 的逐次\Term{group-extension}{扩张}，其中 $\ell_i$ 与 $p$ 互素，且 $\mathbf Z/\ell_i\simeq\mu_{\ell_i}$。
\SourcePage{321}
已知 $\operatorname{Lie}(\operatorname{Pic}_X)=H^1(X,\mathscr O_X)=0$。由 1.2.3 及\Term{section}{截面}的存在，所考察的\Term{functor}{函子} $f^*$ \Term{fully-faithful}{完全忠实}；故只须证明 $\mathrm R^1f_*(f^*G)=0$，或者经逐层\Term{resolution}{分解}，证明
\[\mathrm R^1f_*\mathbf G_a=\mathrm R^1f_*\alpha_p=\mathrm R^1f_*\mu_p=\mathrm R^1f_*\mathbf Z/p=\mathrm R^1f_*\mathbf Z/\ell_i=0.\]
对于 $\mathbf G_a$，这来自 $H^1(X,\mathscr O_X)=0$。$\alpha_p$（相应地 $\mathbf Z/p$）的情形，利用 1.2.3，由以下\Term{exact-sequence}{正合序列}的\Term{cohomology}{上同调}长\Term{exact-sequence}{正合序列}得到：
\[0\longrightarrow\alpha_p\longrightarrow\mathbf G_a\xrightarrow{x^p}\mathbf G_a\longrightarrow0\]
（相应地，$p\ne1$ 时，$0\to\mathbf Z/p\to\mathbf G_a\xrightarrow{x-x^p}\mathbf G_a\to0$）。
$\mathbf Z/\ell_i$（相应地 $\mu_p$）的情形，由\Term{isomorphism}{同构} $\mathbf Z/\ell_i\simeq\mu_{\ell_i}$、1.2.3 以及 Kummer \Term{exact-sequence}{正合序列}
\[0\longrightarrow\mu_n\longrightarrow\mathbf G_m\xrightarrow{x^n}\mathbf G_m\longrightarrow0\]
的长\Term{exact-sequence}{正合序列}得到。
\end{proof}
\begin{SGAresult}{引理}{1.2.9}
设 $S$ 为\Term{local-artinian}{局部阿廷}\Term{scheme}{概形}，$s$ 为 $S$ 的\Term{geometric-point}{几何点}，$f:X\to S$ 为\Term{proper}{固有}\Term{flat}{平坦}\Term{morphism}{态射}，$G$ 为 $S$ 上\Term{flat}{平坦}\Term{finite-type}{有限型}的交换\Term{group-scheme}{群概形}。假设 $f$ \Term{smooth}{光滑}或 $G$ \Term{affine}{仿射}。若 $H^0(X_s,\mathscr O)=k(s)$，$\operatorname{Pic}^\tau_{X_s}=0$，且 $G_s$ 不含\Term{subgroup}{子群} $\mathbf G_m$，则从 $S$ 上的 $G$-\Term{torsor}{挠子}\Term{category}{范畴}到 $X$ 上的 $f^*G$-\Term{torsor}{挠子}\Term{category}{范畴}的\Term{functor}{函子} $f^*$ 是\Term{equivalence}{等价}。
\end{SGAresult}
\begin{proof}
a. 以\Term{finite}{有限}\Term{flat}{平坦}\Term{morphism}{态射} $u:S'\to S$ 作为\Term{localization-morphism}{局部化态射}时，此问题在 $S$ 上是\Term{local-scheme}{局部}的。因此可以假设 $f$ 有\Term{section}{截面} $x$。由 1.2.3 及 $x$ 的存在，\Term{functor}{函子} $f^*$ \Term{fully-faithful}{完全忠实}；只须证明每个沿 $x$ \Term{trivialization}{平凡化}的\Term{torsor}{挠子}均平凡。
b. 假设 $S$ 为域的谱。若 $G$ \Term{affine}{仿射}，断言来自 1.2.8。若 $X$ \Term{smooth}{光滑}，且 $G$ 是\Term{abelian-variety}{阿贝尔簇} $A$ 被\Term{affine}{仿射}群 $G_0$ 的\Term{group-extension}{扩张}，以 $G^n$ 表示 ${}_nA$ 在 $G$ 中的\Term{inverse-image}{逆像}。群 $H^1(X,A)$ 是\Term{torsion-group}{挠群}（Raynaud [1, XIII 2.6]），故为各 $H^1(X,{}_nA)$ 的像之并。因此 $H^1(X,G)$ 为各 $H^1(X,G^n)$ 的像之并。由于 $G^n$ \Term{affine}{仿射}，再由 1.2.8 得到结论。
c. 在假定存在\Term{section}{截面} $x$ 的前提下，对 $S$ 的长度归纳，证明一般情形。设 $i:S_0\hookrightarrow S$ 为由平方为零且长度为 $1$ 的\Term{ideal}{理想} $I$ 定义的子\Term{scheme}{概形}。
由归纳假设，$X$ 上使 $x^*P$ 平凡的 $G$-\Term{torsor}{挠子} $P$，是 $X\times S_0$ 上平凡 $G$-\Term{torsor}{挠子} $P_0$ 的\Term{deformation}{形变}。若 $G$ \Term{smooth}{光滑}，且 $L$ 为 $G_s$ 的 \Term{lie-algebra}{Lie 代数}，则这些\Term{deformation}{形变}由
\[H^1(X_s,f^*(L\otimes I))=0\]
分类，这就证明了此时的 1.2.9。
对 $S$ 上\Term{flat}{平坦}的 $G$，此论证可用 $G_s$ 的 Lie \Term{complex}{复形} $L$ 推广。若 $L\Omega(G_s/k(s))$ 是 $G_s$ 的相对\Term{cotangent-complex}{余切复形}，则 $L\in\mathrm D^b(k(s))$ 是 $Le^*(L\Omega(G_s/k(s)))$ 的\Term{duality}{对偶}。有 $H^i(L)=0$，当 $i\ne0,1$。
\SourcePage{322}
$S_0$ 上（相应地 $X_0$ 上）平凡 $G$-\Term{torsor}{挠子}到 $S$（相应地 $X$）的\Term{deformation}{形变}由 $H^1(L\otimes I)$（相应地 $H^1(X_s,f^*(L\otimes I))$）分类（Illusie [1]）。由假设，$H^1(X_s,H^0(f^*(L\otimes I)))=0$（因 $H^1(X_s,\mathscr O)=0$），且
\[H^1(L\otimes I)\simeq H^0(X_s,H^1(f^*L\otimes I)).\]
由此推出
\[H^1(L\otimes I)\xrightarrow\sim H^1(X_s,f^*(L\otimes I)),\]
即 $X_0$ 上平凡 $G$-\Term{torsor}{挠子}到 $X$ 的每个\Term{deformation}{形变}，恰为 $S_0$ 上平凡 $G$-\Term{torsor}{挠子}到 $S$ 的唯一\Term{deformation}{形变}的\Term{inverse-image}{逆像}。因此，若此\Term{deformation}{形变}沿某\Term{section}{截面}平凡，则它平凡，从而证明 1.2.9。
\end{proof}
\textbf{1.2.10.} 当 $S$ 阿廷时证明 1.2.2。设 $T$ 为 $G$ 的最大子\Term{torus}{环面}；由于已有\Term{fully-faithful}{完全忠实}性（1.2.4），问题对 $S$ 的\Term{etale-topology}{平展拓扑}是\Term{local-scheme}{局部}的，可假设 $T$ \Term{diagonalizable}{可对角化}：$T\simeq\mathbf G_m^n$。由 1.2.3、1.2.9，\Term{cohomology}{上同调}\Term{exact-sequence}{正合序列}给出
\[\mathrm R^1f_*T\xrightarrow\sim\mathrm R^1f_*G.\]
断言由此得到，因为它对 $G=T$ 显然，且 $\operatorname{Hom}(\mathbf G_m,T)\xrightarrow\sim\operatorname{Hom}(\mathbf G_m,G)$（参见 1.2.2.1）。
\textbf{1.2.11.} 当 $S$ 为\Term{complete}{完备}\Term{noetherian}{诺特}\Term{local-scheme}{局部}\Term{scheme}{概形}时证明 1.2.2。设 $s$ 为 $p$ 的一个\Term{section}{截面}。若 $K$ 是 $\mathbf P(\mathscr E)$ 上的 $G$-\Term{torsor}{挠子}，则由 1.2.10，存在唯一的二元组 $(\widehat\varphi,\widehat\Psi)$，其中 $\widehat\varphi:\mathbf G_m\to G$ 为\Term{form}{形式}\Term{homomorphism}{同态}，而
\[\widehat\Psi:p^*s^*K+\widehat\varphi\mathscr O(1)\xrightarrow\sim\widehat K\]
为\Term{form}{形式}\Term{isomorphism}{同构}。若 $G$ \Term{affine}{仿射}，则 $\widehat\varphi$ \Term{algebraizable}{可代数化}（SGA 3 IX 7.1）。暂且承认 $\widehat\varphi$ 总\Term{algebraizable}{可代数化}为 $\varphi:\mathbf G_m\to G$，我们将在 1.2.13 中证明这一点。
若 $K_0=s^*K$，$K_1=p^*K_0+\varphi\mathscr O(1)$，则有\Term{form}{形式}\Term{isomorphism}{同构}
\[\widehat\Psi:\widehat K_1=K(K_0,\varphi)^\wedge\xrightarrow\sim\widehat K,\]
即 $K-K_1$ 的\Term{form}{形式}\Term{trivialization}{平凡化}。由\Term{form}{形式} GAGA，这个\Term{section}{截面}是\Term{algebra}{代数}的，从而证明 1.2.11。更确切地，若 $G$ \Term{quasi-affine}{拟仿射}，则 $K-K_1$ \Term{representable}{可表}，并应用 EGA III 5.1.4。在一般情形，对\Term{algebraic-space}{代数空间} $K-K_1$ 作同样处理。
\textbf{1.2.12. $\widehat\varphi$ 的\Term{algebra}{代数}性：底为\Term{trait}{离散赋值环谱}时的预备核对。} 假设 $S$ 是\Term{complete}{完备}\Term{trait}{离散赋值环谱}，$S=\operatorname{Spec}(V)$，$L$ 为 $V$ 的\Term{field-of-fractions}{分式域}。
a) 在 $L$ 上存在唯一三元组 $(K_{0L},\varphi_L,\Psi_L)$：
\[\Psi_L:K(K_{0L},\varphi_L)\xrightarrow\sim K|_{\operatorname{Spec}(L)}.\]
b) 由 (1.2.10)、(1.2.11)，采用 1.2.7 的记号，有\Term{rigidity}{刚性}解析映射 $\widehat\varphi_L:U_m\to\widehat G$，一个 $\widehat G$-\Term{torsor}{挠子} $\widehat K_0$，以及\Term{rigidity}{刚性}解析\Term{isomorphism}{同构}
\[\widehat\Psi:p^*\widehat K_0+\widehat\varphi_L\mathscr O(1)^\wedge\xrightarrow\sim\widehat K.\]
由唯一性断言 1.2.7，$\widehat\varphi_L$ 与由 $\varphi_L$ \Term{derive}{导出}的映射相同。因此 1.2.11 的\Term{form}{形式}\Term{morphism}{态射} $\widehat\varphi$ 在 $\mu_n$ 上诱导的映射，与\Term{algebra}{代数}映射 $\varphi_L:\mathbf G_m\to G$ \Term{derive}{导出}的映射一致。
\SourcePage{323}
\textbf{1.2.13. $\widehat\varphi$ 的\Term{algebra}{代数}性。} 1.2.11 的\Term{form}{形式}映射 $\widehat\varphi:\mathbf G_m\to G$ 对 $S$ 上\Term{invertible}{可逆}的 $\ell$ 诱导映射
\[\varphi_\ell:T_\ell(\widehat{\mathbf G}_m)=T_\ell(\mathbf G_m)\longrightarrow T_\ell(G).\]
由 (1.2.12) 及 EGA II 7.1.9，对 $S$ 的每点 $t$，都有一个\Term{algebra}{代数}\Term{morphism}{态射} $\varphi_t:\mathbf G_{mt}\to G_t$ 满足 $T_\ell(\varphi_t)=T_\ell(\varphi)_t$。故由 SGA 3 XV 4.1 bis，$\widehat\varphi$ \Term{algebraizable}{可代数化}。
\textbf{1.2.14.} 证明 1.2.2。取\Term{limit}{极限}，可归结于 $S$ \Term{noetherian}{诺特}，然后归结于 $S$ 为\Term{affine}{仿射}\Term{algebra}{代数} $A$ 的\Term{noetherian}{诺特}\Term{local-scheme}{局部}\Term{scheme}{概形}。设 $s$ 为 $p$ 的一个\Term{section}{截面}，$K$ 为 $\mathbf P(\mathscr E)$ 上的 $G$-\Term{torsor}{挠子}。在 $S'=\operatorname{Spec}(\widehat A)$ 上，存在唯一\Term{homomorphism}{同态} $\varphi':\mathbf G_{mS'}\to G_{S'}$，使 $K$ 在 $\mathbf P(\mathscr E)\times_S S'$ 上的\Term{inverse-image}{逆像} $K'$ 具有\Term{form}{形式} $K(K'_0,\varphi')$。\Term{homomorphism}{同态} $\varphi'$ \Term{descent}{下降}为 $\varphi:\mathbf G_m\to G$。为证明\Term{functor}{函子} 1.2.1 的\Term{essentially-surjective}{本质满射}性，可将 $K$ 换为 $K-\varphi\mathscr O(1)$，即假设 $\varphi=0$。在 $S'$ 上，此时存在唯一\Term{isomorphism}{同构} $\Psi:p'^*s^*K'\xrightarrow\sim K'$，使它在\Term{section}{截面} $s$ 上诱导恒等。这个\Term{morphism}{态射}\Term{descent}{下降}至 $S$（fppf \Term{torsor}{挠子}\Term{morphism}{态射}的 fpqc \Term{descent}{下降}），从而证明 1.2.2。
\begin{SGAresult}{变式}{1.2.15}
1.2.2 的结论在下列情形仍成立：
(i) $G$ 由 $S$ 上的\Term{quasi-coherent}{拟凝聚}模定义；
(ii) $G$ 为 $S$ 的小\Term{etale-site}{平展景}上\Term{torsion-sheaf}{挠层}的\Term{inverse-image}{逆像}。
\end{SGAresult}
在情形 (i)，由\Term{proper}{固有}\Term{morphism}{态射}的\Term{base-change}{基变换}定理，以及 $k$ \Term{algebraically-closed}{代数闭}时 $\mathbf P^r_k$ 的\Term{simply-connected}{单连通}性，\Term{arrow}{箭头} (1.2.2.1) 为\Term{isomorphism}{同构}。
在情形 (ii)，鉴于 $G$ 的 fppf \Term{cohomology}{上同调}与它的 Zariski \Term{topology}{拓扑}\Term{cohomology}{上同调}一致，\Term{arrow}{箭头} (1.2.2.1) 为\Term{isomorphism}{同构}是标准结果。
\SGAsection{1.3}{挠子的积分}
\textbf{1.3.1.} 设 $S$ 为\Term{scheme}{概形}，$G$ 为 $S$ 的大 fppf 景上的\Term{abelian-sheaf}{阿贝尔层}，满足以下三个条件。
\textbf{1.3.1.1.}
(i) 对每个\Term{projective-space}{射影空间} $p:\mathbf P^r_S\to S$，$r\geq1$，在 $S$ 的大 fppf 景（XVII 0.10）中有
\[G\xrightarrow\sim p_*p^*G,\qquad\mathscr Hom(\mathbf G_m,G)\xrightarrow\sim\mathrm R^1p_*p^*G.\]
(ii) 对每个 $n$ 及每个\Term{affine}{仿射}\Term{smooth}{光滑}的 $S$-\Term{morphism}{态射} $f:X\to T$，若 $K_1,K_2$ 为 $\operatorname{Sym}^n_T(X)$ 上的两个 $G$-\Term{torsor}{挠子}，则\Term{set}{集合} $\operatorname{Hom}(K_1,K_2)$ 与 $K_1$ 在 $(X/T)^n$ 上的\Term{inverse-image}{逆像}到 $K_2$ 在 $(X/T)^n$ 上的\Term{inverse-image}{逆像}的对称\Term{morphism}{态射}集相同。
\SourcePage{324}
此条件\Term{equivalence}{等价}于以下两者同时成立：
\[(\mathrm{ii}')\quad\operatorname{Hom}_S(\operatorname{Sym}^n_T(X),G)\xrightarrow\sim\operatorname{Hom}_S((X/T)^n,G)^{\mathfrak S_n};\]
$(\mathrm{ii}'')$ 若 $\operatorname{Sym}^n_T(X)$ 上\Term{torsor}{挠子} $K$ 在 $(X/T)^n$ 上的\Term{inverse-image}{逆像}有一个对称\Term{section}{截面}，则 $K$ 平凡。
(iii) $G$ 满足条件 XVII 6.3.21.1，或同时满足条件 XVII 6.3.24.1、XVII 6.3.24.2。
因此 1.2.2 的结论成立（参见 1.2.2.1），公式汇编 XVII 6.3.26 也成立。
条件 (i) 已在 1.2.2、1.2.15 中讨论。当 $G$ 为\Term{finitely-presented}{有限表示}的\Term{representable}{可表}\Term{sheaf}{层}，或是 $S$ 的小\Term{etale-site}{平展景}上某\Term{sheaf}{层}的\Term{inverse-image}{逆像}，或为\Term{flat}{平坦}\Term{quasi-coherent}{拟凝聚}\Term{sheaf}{层}时，条件 (ii) 均成立。故条件 1.3.1.1 在下列每个情形中成立。
\textbf{1.3.1.2.}
(a) $G$ \Term{smooth}{光滑}且\Term{finitely-presented}{有限表示}（1.2.2、XVII 6.3.21.1）；
(b) $G$ 为 $S$ 的小\Term{etale-site}{平展景}上某\Term{torsion-sheaf}{挠层}的\Term{inverse-image}{逆像}（1.2.15 (ii)、XVII 6.3.21.1）；
(c) $G$ \Term{affine}{仿射}，且为\Term{smooth}{光滑}群之间一个满\Term{homomorphism}{同态}的核（1.2.2、XVII 6.3.3.1）；
(d) $G$ 由 $S$ 上某\Term{flat}{平坦}\Term{quasi-coherent}{拟凝聚}\Term{sheaf}{层}定义（1.2.15 (i)、XVII 6.3.21.1）。
\textbf{1.3.2.} 设 $f:X\to S$ 为\Term{scheme}{概形} $S$ 上的\Term{projective}{射影}\Term{smooth}{光滑}曲线，其\Term{geometric-fiber}{几何纤维}\Term{connected-nonempty}{连通非空}。设 $\mathscr L$ 为 $X$ 上的\Term{invertible-sheaf}{可逆层}，并在以下意义下“充分”相对\Term{ample}{丰沛}。
\textbf{1.3.2.1.} 对 $S$ 的每个\Term{geometric-point}{几何点} $s$，若 $X_s$ 的\Term{genus}{亏格}为 $g$，则
\[\deg_{X_s}(\mathscr L)>2g-2,\qquad\deg_{X_s}(\mathscr L)\geq1\ (g=0),\qquad\deg_{X_s}(\mathscr L)\geq2\ (g=1).\]
于是 $\mathrm R^1f_*\mathscr L=0$，而 $f_*\mathscr L$ \Term{locally-free}{局部自由}，其构造与任意\Term{base-change}{基变换}相容，且处处秩 $\geq2$。若 $f_*\mathscr L$ 的\Term{section}{截面} $s$ 处处不为零，则将 $s$ 视作 $X$ 上 $\mathscr L$ 的\Term{section}{截面}时，其零点\Term{scheme}{概形} $D_s$ 是一个\Term{relative-divisor}{相对除子}（EGA IV 21.15.2）；由于 $X/S$ 的\Term{relative-dimension}{相对维数}为 $1$，它在 $S$ 上\Term{finite}{有限}\Term{locally-free}{局部自由}。\Term{canonical}{典范}地有
\[s:\mathscr O(D_s)\xrightarrow\sim\mathscr L.\]
\Term{divisor}{除子} $D_s$ 只依赖于 $s$ 模乘以 $\mathscr O_S^*$ 的一个\Term{section}{截面}。因此上述构造在由\Term{base-change}{基变换} $p:\mathbf P(f_*(\mathscr L)^\vee)\to S$ 从 $X$ 得到的\Term{scheme}{概形}上给出一个\Term{relative-divisor}{相对除子} $Z_{\mathscr L}$。
设 $K$ 为 $G_X$ 下的\Term{torsor}{挠子}，$K|Z_{\mathscr L}$ 为它在 $Z_{\mathscr L}$ 上的\Term{inverse-image}{逆像}。这个\Term{torsor}{挠子}从 $Z_{\mathscr L}$ 到 $\mathbf P(f_*(\mathscr L)^\vee)$ 的迹（XVII 6.3.26）是 $\mathbf P(f_*(\mathscr L)^\vee)$ 上的\Term{torsor}{挠子} $K_{\mathscr L}$：
\[K_{\mathscr L}=\operatorname{Tr}_{Z_{\mathscr L}/\mathbf P}(K|Z_{\mathscr L}).\]
\SourcePage{325}
由 1.2.2，存在一个在唯一\Term{isomorphism}{同构}意义下唯一的三元组，由 $S$ 上的 $G$-\Term{torsor}{挠子} $\langle\mathscr L,K]$、\Term{morphism}{态射} $\varphi_{\mathscr L,K}:\mathbf G_m\to G$ 及\Term{isomorphism}{同构}
\[K_{\mathscr L}\simeq p^*\langle\mathscr L,K]+\varphi_{\mathscr L,K}\mathscr O(1)\]
组成。
\textbf{1.3.3.} 若将\Term{torsor}{挠子}的合成写作加法，则\Term{canonical}{典范}地有
\begin{align}\langle\mathscr L,K_1+K_2]&\simeq\langle\mathscr L,K_1]+\langle\mathscr L,K_2],\tag{1.3.3.1}\\\varphi_{\mathscr L,K_1+K_2}&=\varphi_{\mathscr L,K_1}+\varphi_{\mathscr L,K_2}.\tag{1.3.3.2}\end{align}
设 $\mathscr L_1,\mathscr L_2$ 为 $X$ 上满足 (1.3.3.1) 的两个\Term{invertible-sheaf}{可逆层}，$\mathscr L=\mathscr L_1\otimes\mathscr L_2$。\Term{section}{截面}的\Term{tensor-product}{张量积}定义\Term{morphism}{态射}
\[\begin{tikzcd}[column sep=small]\mathbf P((f_*\mathscr L_1)^\vee)\times_S\mathbf P((f_*\mathscr L_2)^\vee)\ar[rr,"\pi"]\ar[d,"\operatorname{pr}_1"']\ar[dr,"p_{12}" description]\ar[rrd,"\operatorname{pr}_2" near start]&&\mathbf P((f_*\mathscr L)^\vee)\ar[dd,"p"]\\\mathbf P((f_*\mathscr L_1)^\vee)\ar[dr,"p_1"']&S&\mathbf P((f_*\mathscr L_2)^\vee)\ar[dl,"p_2"]\\&S\ar[u,equal]\ar[r,equal]&S.\end{tikzcd}\]
有 $\pi^*Z_{\mathscr L}=Z_{\mathscr L_1}+Z_{\mathscr L_2}$，故得到\Term{isomorphism}{同构}（XVII 6.3.27.1）
\begin{equation}\pi^*K_{\mathscr L}=\operatorname{pr}_1^*K_{\mathscr L_1}+\operatorname{pr}_2^*K_{\mathscr L_2}.\tag{1.3.3.3}\end{equation}
有\Term{canonical-isomorphism}{典范同构} $\pi^*\mathscr O(1)\simeq\operatorname{pr}_1^*\mathscr O(1)\otimes\operatorname{pr}_2^*\mathscr O(1)$，所以 (1.3.3.3) 可改写为
\[\begin{aligned}\pi^*(p^*\langle\mathscr L,K]+\varphi_{\mathscr L,K}\mathscr O(1))
&=p_{12}^*\langle\mathscr L,K]+\operatorname{pr}_1^*\varphi_{\mathscr L,K}\mathscr O(1)+\operatorname{pr}_2^*\varphi_{\mathscr L,K}\mathscr O(1)\\
&=p_{12}^*\langle\mathscr L_1,K]+p_{12}^*\langle\mathscr L_2,K]+\operatorname{pr}_1^*\varphi_{\mathscr L_1,K}\mathscr O(1)+\operatorname{pr}_2^*\varphi_{\mathscr L_2,K}\mathscr O(1).
\end{aligned}\]
两次应用 1.2.2，得到\Term{canonical-isomorphism}{典范同构}
\begin{equation}\langle\mathscr L_1\otimes\mathscr L_2,K]=\langle\mathscr L_1,K]+\langle\mathscr L_2,K]\tag{1.3.3.4}\end{equation}
以及恒等式 $\varphi_{\mathscr L_1\otimes\mathscr L_2,K}=\varphi_{\mathscr L_1,K}=\varphi_{\mathscr L_2,K}$。
所以 $\varphi_{\mathscr L,K}$ 只依赖于 $K$；置
\begin{equation}\deg(K)=\varphi_{\mathscr L,K}.\tag{1.3.3.5}\end{equation}
\SourcePage{326}
于是有
\begin{align}K_{\mathscr L}&\simeq p^*\langle\mathscr L,K]+\deg(K)\mathscr O(1),\tag{1.3.3.6}\\\deg(K_1+K_2)&=\deg(K_1)+\deg(K_2).\tag{1.3.3.7}\end{align}
设 $\mathscr U$ 为 $f_*\mathscr L$ 的一个\Term{local-scheme}{局部}\Term{direct-summand}{直和因子}\Term{invertible-sheaf}{可逆层}，定义 $\mathbf P((f_*\mathscr L)^\vee)$ 的\Term{section}{截面} $u$。以 $D_u$ 表示 $\mathscr U$ 的\Term{local-scheme}{局部}\Term{invertible}{可逆}\Term{section}{截面}定义的\Term{relative-divisor}{相对除子}。有 $u^*\mathscr O(1)\simeq\mathscr U^\vee$，故由 (1.3.3.6) 得到
\begin{equation}u^*K_{\mathscr L}=\operatorname{Tr}_{D/S}(K)\simeq\langle\mathscr L,K]-\deg(K)\mathscr U.\tag{1.3.3.8}\end{equation}
以下相容性留给读者核对。
\begin{SGAresult}{命题}{1.3.4}
由\Term{isomorphism}{同构} (1.3.3.1)、(1.3.3.4) 组成的下图交换。为使图紧凑，写 $A_{ij}=\langle\mathscr L_i,K_j]$：
\[\begin{tikzcd}[column sep=small]\langle\mathscr L_1\otimes\mathscr L_2,K_1+K_2]\ar[r,"\sim"]\ar[d,"\sim"']&\langle\mathscr L_1\otimes\mathscr L_2,K_1]+\langle\mathscr L_1\otimes\mathscr L_2,K_2]\ar[d,"\sim"]\\\langle\mathscr L_1,K_1+K_2]+\langle\mathscr L_2,K_1+K_2]\ar[d,"\sim"']&A_{11}+A_{21}+A_{12}+A_{22}\ar[d,"\sim"]\\A_{11}+A_{12}+A_{21}+A_{22}\ar[r,equal]&A_{11}+A_{12}+A_{21}+A_{22}.\end{tikzcd}\]
\end{SGAresult}
\Term{isomorphism}{同构} (1.3.3.1)、(1.3.3.4) 还与结合及交换\Term{isomorphism}{同构}相容，其确切意义留给读者阐明。于是公式 (1.3.3.4) 允许按“线性”延拓 $\langle\mathscr L,K]$ 的定义，使 $\mathscr L$ 不必再满足条件 (1.3.2.1)。
\textbf{1.3.5.} 对 $X$ 上的每个\Term{relative-divisor}{相对除子} $D$，置
\begin{equation}\operatorname{Tr}_{D/S}(K)=\operatorname{Tr}_{D/S}(K|D).\tag{1.3.5.0}\end{equation}
若 $\mathscr O(D)$ 满足 (1.3.2.1)，则 $\mathscr O(D)$ 的\Term{section}{截面} $1$ 为 $f_*\mathscr O(D)$ 的\Term{section}{截面}，并在后者中生成一个\Term{local-scheme}{局部}为\Term{direct-summand}{直和因子}的\Term{invertible-sheaf}{可逆层}，且给出该\Term{sheaf}{层}的\Term{trivialization}{平凡化}。因此\Term{isomorphism}{同构} (1.3.3.8) 给出\Term{isomorphism}{同构}
\begin{equation}\langle\mathscr O(D),K]\simeq\operatorname{Tr}_{D/S}(K).\tag{1.3.5.1}\end{equation}
\SourcePage{327}
立即验证可知，以下两图交换：
\begin{equation}\begin{tikzcd}\langle\mathscr O(D),K_1+K_2]\ar[r,"\sim"]\ar[d,"\sim"']&\operatorname{Tr}_{D/S}(K_1+K_2)\ar[d,"\sim"]\\\langle\mathscr O(D),K_1]+\langle\mathscr O(D),K_2]\ar[r,"\sim"]&\operatorname{Tr}_{D/S}(K_1)+\operatorname{Tr}_{D/S}(K_2)\end{tikzcd}\tag{1.3.5.2}\end{equation}
及
\begin{equation}\begin{tikzcd}\langle\mathscr O(D_1+D_2),K]\ar[r,"\sim"]\ar[d,"\sim"']&\operatorname{Tr}_{D_1+D_2/S}(K)\ar[d,"{\sim\ (\text{XVII }6.3.27.1)}"]\\\langle\mathscr O(D_1),K]+\langle\mathscr O(D_2),K]\ar[r,"\sim"]&\operatorname{Tr}_{D_1/S}(K)+\operatorname{Tr}_{D_2/S}(K).\end{tikzcd}\tag{1.3.5.3}\end{equation}
因而可按线性，对任意\Term{relative-divisor}{相对除子}定义\Term{isomorphism}{同构} (1.3.5.1)。
\textbf{1.3.6.} 设 $D,E$ 为两个\Term{effective}{有效}相对 \Term{cartier-divisor}{Cartier 除子}，$f$ 为满足 $\operatorname{div}(f)=D-E$ 的\Term{rational-function}{有理函数}，即 $\mathscr O(D)$ 与 $\mathscr O(E)$ 之间的一个\Term{isomorphism}{同构}。$\mathscr O(D)$ 到 $\mathscr O(E)$ 的\Term{isomorphism}{同构} $f$，由 (1.3.5.1) 定义\Term{isomorphism}{同构}
\begin{equation}\langle f,K]:\operatorname{Tr}_{D/S}(K)\xrightarrow\sim\operatorname{Tr}_{E/S}(K).\tag{1.3.6.1}\end{equation}
有
\[\langle fg,K]=\langle f,K]\langle g,K],\qquad\langle f,K_1+K_2]=\langle f,K_1]+\langle f,K_2],\]
其中第二式来自 (1.3.5.2)。
若 $\lambda$ 为 $\mathscr O_S^*$ 的\Term{section}{截面}，则对 $X$ 上任意\Term{invertible-sheaf}{可逆层} $\mathscr L$，乘以 $\lambda$ 是 $\mathscr L$ 的\Term{automorphism}{自同构}。搬运结构，此\Term{automorphism}{自同构}在 (1.3.3.8) 两侧诱导彼此由该\Term{isomorphism}{同构}对应的\Term{automorphism}{自同构}。因此得到
\[\langle\lambda,K]=\deg(K)(\lambda).\]
若 $\deg(K)\ne0$，则 $\operatorname{Tr}_{D/S}(K)$ 与 $\operatorname{Tr}_{E/S}(K)$ 之间的\Term{isomorphism}{同构} $\langle f,K]$ 依赖于 $f$ 的选取。
设 $\lambda\in\Gamma(S,\mathscr O_S^*)$，$g\in\Gamma(S,G)$，$\mathscr L$ 为 $X$ 上的\Term{invertible-sheaf}{可逆层}，$K$ 为 $X$ 上的 $G$-\Term{torsor}{挠子}。若 $\langle\lambda,g]$ 为由 $\mathscr L$ 的\Term{automorphism}{自同构} $\lambda$ 和 $K$ 的\Term{automorphism}{自同构} $g$ 搬运结构而得的 $\langle\mathscr L,K]$ 的\Term{automorphism}{自同构}，则将它视为 $G$ 的一个\Term{section}{截面}时，有 $\langle\lambda,g]=\langle\lambda,K]+\langle\mathscr L,g]$；由于 $\langle\mathscr O(D),g]=\operatorname{Tr}_{D/S}(g)$，从上式得到
\begin{equation}\langle\lambda,g]=\deg(K)(\lambda)+\deg(\mathscr L)g.\tag{1.3.6.2}\end{equation}
\SourcePage{328}
\textbf{1.3.7.} 设 $f:X\to S$ 为 $S$ 上的\Term{projective}{射影}\Term{smooth}{光滑}曲线，$G$ 为 $S$ 上满足 1.3.1.1 的群。若 $X\xrightarrow{f'}T\xrightarrow\Pi S$ 为 $f$ 的 \Term{stein-factorization}{Stein 分解}，则 $\Pi$ 为 $S$ 的\Term{etale-cover}{平展覆盖}，\Term{sheaf}{层} $f_*\mathbf G_m$ 为\Term{torus}{环面} $\prod_{T/S}\mathbf G_m$。在 $S$ 上对\Term{etale-topology}{平展拓扑}\Term{local-scheme}{局部}地，$X$ 是一族 $S$ 上\Term{geometric-fiber}{几何纤维}\Term{connected-nonempty}{连通非空}的曲线 $(X_i)_{i\in I}$ 之和。若 $K$ 为 $G$ 下的\Term{torsor}{挠子}，$\mathscr L$ 为 $X$ 上的\Term{invertible-sheaf}{可逆层}，则置
\[\langle\mathscr L,K]=\sum_i\langle\mathscr L|X_i,K|X_i],\]
并定义 $K$ 的次数 $f_*\mathbf G_m=\mathbf G_m^I\to G$，使其坐标为 $\deg(K|X_i):\mathbf G_m\to G$。这些\Term{local-scheme}{局部}构造可整体化，给出\Term{functor}{函子}和\Term{morphism}{态射}
\begin{align}\langle\mathscr L,K]_{X/S}&=\operatorname{Tr}_{T/S}(\langle\mathscr L,K]_{X/T}),\tag{1.3.7.1}\\\deg(K):f_*\mathbf G_m=\prod_{T/S}\mathbf G_m&\longrightarrow G,\qquad\deg_{X/S}(K)=\prod_{T/S}\deg_{X/T}(K).\tag{1.3.7.2}\end{align}
对由 $\Gamma(T,\mathscr O_T^*)$ 中的元素 $\lambda_0$ 定义的 $\mathscr L$ 的\Term{automorphism}{自同构} $\lambda$，有
\begin{equation}\langle\lambda,K]=\deg(K)(\lambda).\tag{1.3.7.3}\end{equation}
将 $K$ 的总次数定义为复合
\[\deg_{\mathrm{tot}}(K):\mathbf G_m\longrightarrow f_*\mathbf G_m\xrightarrow{\deg(K)}G.\]
以上结果给出以下公式汇编。
\begin{SGAresult}{公式汇编}{1.3.8}
对每条\Term{projective}{射影}\Term{smooth}{光滑}曲线 $f:X\to S$，有构造与任意\Term{base-change}{基变换}相容的\Term{functor}{函子} $\langle\ ,\ ]$，它将 $X$ 上的\Term{invertible-sheaf}{可逆层} $\mathscr L$ 和 $f^*G$ 下的\Term{torsor}{挠子} $K$ 送到 $S$ 上的 $G$-\Term{torsor}{挠子} $\langle\mathscr L,K]$。此\Term{functor}{函子}带有下列附加结构，并满足下列条件。
\end{SGAresult}
\textbf{1.3.8.1.} $\langle\mathscr L,K]$ 对 $\mathscr L,K$ 双\Term{additive}{加性}；双\Term{additive}{加性}\Term{isomorphism}{同构}与结合、交换\Term{isomorphism}{同构}以及\Term{base-change}{基变换}相容，满足相容性 (1.3.4)，并且通过 XVII 6.3.25.3、XVII 6.3.27.1 与下面的\Term{isomorphism}{同构} (1.3.8.2) 相容。
\textbf{1.3.8.2.} 对 $X/S$ 上每个\Term{effective}{有效}\Term{relative-divisor}{相对除子} $E$，有\Term{canonical-isomorphism}{典范同构}
\[\langle\mathscr O(E),K]\simeq\operatorname{Tr}_{E/S}(K|E).\]
\SourcePage{329}
\textbf{1.3.8.3.} 次数 (1.3.7.2) 对 $\mathscr L$ 的\Term{automorphism}{自同构} $\lambda\in f_*\mathbf G_m$ 满足 $\langle\lambda,K]=\deg(K)(\lambda)$。
对 $\lambda\in\mathbf G_m(S)$，总次数满足 $\langle\lambda,K]=\deg_{\mathrm{tot}}(K)(\lambda)$。
同样，对 $g\in G(S)$，有 $\langle\mathscr L,g]=\deg(\mathscr L)g$。
由此得到\Term{canonical-isomorphism}{典范同构}
\begin{equation}\langle f^*\mathscr L,K]\simeq\deg_{\mathrm{tot}}(K)(\mathscr L).\tag{1.3.8.4}\end{equation}
由定义立即得到：对每个\Term{homomorphism}{同态} $\varphi:G\to G'$，有与上述数据相容的\Term{canonical-isomorphism}{典范同构}
\begin{equation}\langle\mathscr L,\varphi(K)]\simeq\varphi(\langle\mathscr L,K]).\tag{1.3.8.5}\end{equation}
\textbf{1.3.8.6.} 设 $u:X\to Y$ 为 $S$ 上\Term{projective}{射影}\Term{smooth}{光滑}曲线之间的\Term{flat}{平坦}\Term{morphism}{态射}。则有与双\Term{additive}{加性}及\Term{base-change}{基变换}相容的\Term{isomorphism}{同构}：
对 $X$ 上的\Term{invertible-sheaf}{可逆层} $\mathscr L$ 和 $Y$ 上的 $G_Y$-\Term{torsor}{挠子} $K$，有
\[\langle\mathscr L,u^*K]\xrightarrow\sim\langle N_{X/Y}(\mathscr L),K].\]
对 $X$ 上的\Term{invertible-sheaf}{可逆层} $\mathscr L$ 和 $X$ 上的 $G_X$-\Term{torsor}{挠子} $K$，有
\[\langle u^*\mathscr L,K]\xrightarrow\sim\langle\mathscr L,\operatorname{Tr}_{X/Y}K].\]
此外，这些\Term{isomorphism}{同构}对曲线\Term{morphism}{态射}的复合 $uv$ 满足显然的相容性；对 $\mathscr L=\mathscr O(E)$，它们等同于\Term{isomorphism}{同构} XVII 6.3.27.2：
\[\operatorname{Tr}_{E/S}(u^*K)\xrightarrow\sim\operatorname{Tr}_{p_*E/S}(K),\qquad\operatorname{Tr}_{p_*E/S}(K)\xrightarrow\sim\operatorname{Tr}_{E/S}\operatorname{Tr}_{X/Y}K.\]
尚须证明 1.3.8.6。归结于 $X,Y$ 的\Term{geometric-fiber}{几何纤维}\Term{connected-nonempty}{连通非空}的情形，并注意到上述显式公式对 $\mathscr L$ 的每个\Term{section}{截面} $s$，只要它作为 $\mathscr L$ 的\Term{direct-image}{直像}的\Term{section}{截面}不消失，就给出所需类型的\Term{isomorphism}{同构} $\Psi_s$。它关于 $s$ \Term{homogeneous}{齐次}，因此当 $\mathscr L$ 满足 (1.3.2.1) 时不依赖于 $s$。一般情形由此推出。
\textbf{1.3.9.} 在 1.3.7 的假设下，我们已看到（1.3.8），$X$ 上的每个 $G$-\Term{torsor}{挠子} $K$ 定义一个\Term{functor}{函子} $\langle *,K]$，以下记为 $F(*)$，带有下列附加数据，并满足下列条件：
(i) $F$ 为 $S$ 上的 fppf \Term{stack}{叠}\Term{morphism}{态射}（XVII 0.10、Giraud [1]），
\SourcePage{330}
其源是对 $S$-\Term{scheme}{概形} $T$ 以 $X_T$ 上的\Term{invertible-sheaf}{可逆层}为 $T$-\Term{object}{对象}的\Term{stack}{叠}，其靶为 $G$-\Term{torsor}{挠子}\Term{stack}{叠}。
(ii) \Term{functor}{函子} $F$ 带有\Term{functorial}{函子性}的\Term{additive}{加性}\Term{isomorphism}{同构} $F(\mathscr L_1\otimes\mathscr L_2)\xrightarrow\sim F(\mathscr L_1)+F(\mathscr L_2)$，与结合、交换\Term{isomorphism}{同构}相容，且与\Term{base-change}{基变换}相容。
\begin{SGAresult}{定理}{1.3.10}
在 1.3.7 的假设下，\Term{functor}{函子} $\Phi:K\mapsto\langle *,K]$ 给出 $X$ 上的 $G$-\Term{torsor}{挠子}\Term{category}{范畴}与 1.2.9 中考虑的\Term{stack}{叠}\Term{morphism}{态射}\Term{category}{范畴}之间的\Term{equivalence}{等价}。
\end{SGAresult}
\begin{proof}
容易归结于 $X$ 的\Term{fiber}{纤维}\Term{connected-nonempty}{连通非空}的情形。
设 $F$ 满足 1.3.9 (i)、(ii)。对每个 $S$-\Term{scheme}{概形} $T$ 和 $X_T/T$ 的每个\Term{section}{截面} $t$，仍以 $t$ 表示\Term{relative-divisor}{相对除子} $t(T)$。于是 $F(\mathscr O(t))$ 为 $T$ 上的 $G_T$-\Term{torsor}{挠子}。当 $t$ 为 $X$ 的取值于 $X$ 的“通有点”，即由 $X$ 的恒等映射定义时，以 $\Psi(F)$ 表示如此得到的 $X$ 上的\Term{torsor}{挠子}。对上述每个\Term{section}{截面} $t$，有
\begin{equation}F(\mathscr O(t))\simeq t^*\Psi(F).\tag{1.3.10.1}\end{equation}
\Term{isomorphism}{同构} 1.3.8.2 给出\Term{functor}{函子}\Term{isomorphism}{同构} $\Psi\circ\Phi\simeq\mathrm{Id}$。下面构造\Term{isomorphism}{同构} $\Phi\circ\Psi\simeq\mathrm{Id}$。
对 $X_T/T$ 的每个\Term{section}{截面} $t$，\Term{isomorphism}{同构} (1.3.10.1) 是\Term{isomorphism}{同构}
\[\alpha:F(\mathscr O(t))\xrightarrow\sim\langle\mathscr O(t),\Psi(F)].\]
作加法，$X_T/T$ 的每一族\Term{section}{截面} $t_1,\ldots,t_n$ 定义\Term{isomorphism}{同构}
\[\alpha:F(\mathscr O(\textstyle\sum t_i))\xrightarrow\sim\langle\mathscr O(\textstyle\sum t_i),\Psi(F)],\]
且由 1.3.9 (ii)，此\Term{isomorphism}{同构}不依赖 $t_i$ 的次序。
取 $T=(X/S)^n$，$t_1,\ldots,t_n$ 为通有\Term{section}{截面} $n$ 元组。若 $T'=\operatorname{Sym}^n_S(X)$，则在 $X_{T'}$ 上有\Term{divisor}{除子} $\sum t_i$，使 $X_T$ 上的\Term{divisor}{除子} $\sum t_i$ 为其\Term{inverse-image}{逆像}。在 $T$ 上，$\alpha$ 是两个从 $T'$ 上的\Term{torsor}{挠子}\Term{inverse-image}{逆像}而来的\Term{torsor}{挠子}之间的\Term{symmetry-isomorphism}{对称同构}：
\[\alpha:F(\mathscr O(\textstyle\sum t_i))\xrightarrow\sim\langle\mathscr O(\textstyle\sum t_i),\Psi(F)].\]
由假设 1.3.1.1 (ii)，此\Term{isomorphism}{同构}来自 $T'$ 上唯一的\Term{isomorphism}{同构}
\[\alpha:F(\mathscr O(\textstyle\sum t_i))\xrightarrow\sim\langle\mathscr O(\textstyle\sum t_i),\Psi(F)]\quad(\text{在 }T'\text{ 上}).\]
根据 XVII 6.3，遂对 $X$ 上的每个\Term{relative-divisor}{相对除子} $D$ 得到\Term{isomorphism}{同构}
\begin{equation}\alpha_D:F(\mathscr O(D))\xrightarrow\sim\langle\mathscr O(D),\Psi(F)],\tag{1.3.10.2}\end{equation}
它对 $D$ \Term{additive}{加性}，构造与任意\Term{base-change}{基变换}相容。
\SourcePage{331}
设 $\mathscr L$ 为 $X$ 上在 1.3.2.1 意义下充分\Term{ample}{丰沛}的\Term{invertible-sheaf}{可逆层}。$f_*\mathscr L$ 的每个处处非\Term{zero-section}{零截面} $s$，用 1.3.2 的记号，定义\Term{isomorphism}{同构} $s:\mathscr O(D_s)\xrightarrow\sim\mathscr L$，从而由 (1.3.10.2) 得到\Term{isomorphism}{同构}
\begin{equation}\alpha_s:F(\mathscr L)\xrightarrow\sim\langle\mathscr L,\Psi(F)].\tag{1.3.10.3}\end{equation}
证明这个\Term{isomorphism}{同构}不依赖于 $s$。事实上，(i) $\alpha_s$ 的构造与任意\Term{base-change}{基变换}相容；(ii) 存在\Term{homomorphism}{同态} $\varphi:\mathbf G_m\to G$，使 $\alpha_{\lambda s}=\varphi(\lambda)\alpha_s$（由 (1.3.10.3) 两侧的\Term{functorial}{函子性}）。
必有 $\varphi=0$，否则 $F(\mathscr L)-\langle\mathscr L,\Psi(F)]$ 在 $\mathbf P((f_*\mathscr L)^\vee)$ 上的\Term{inverse-image}{逆像}\Term{torsor}{挠子}将具有非零次数。因此各 $\alpha_s$ 定义了从 $\mathbf P((f_*\mathscr L)^\vee)$ 到 $G$ 的\Term{morphism}{态射}，而由假设 1.3.1.1 (i)，后者为常值。因此可置 $\alpha_{\mathscr L}=\alpha_s$。
这些\Term{isomorphism}{同构} $\alpha_{\mathscr L}$ 组成\Term{functor}{函子} $F$ 和 $\langle *,\Psi(F)]$ 在充分\Term{ample}{丰沛}\Term{sheaf}{层}上的\Term{restriction}{限制}之间的\Term{functor}{函子}\Term{isomorphism}{同构}。这个\Term{isomorphism}{同构}与\Term{additive}{加性}\Term{isomorphism}{同构}相容，所以可以线性地延拓到一切\Term{invertible-sheaf}{可逆层}，完成 $\alpha:\Phi\circ\Psi\simeq\mathrm{Id}$ 的构造。
\end{proof}
\begin{SGAresult}{例}{1.3.11}
$G=\mathbf G_m$。在 1.3.7 的假设下，当 $G=\mathbf G_m$ 时，将 $\langle\ ,\ ]$ 写作 $\langle\ ,\ \rangle$；若 $\mathscr L,\mathscr M$ 为 $X$ 上的两个\Term{invertible-sheaf}{可逆层}，则 $\langle\mathscr L,\mathscr M\rangle$ 为 $S$ 上的\Term{invertible-sheaf}{可逆层}。并且有\Term{isomorphism}{同构} (1.3.8.2)：
\[\langle\mathscr O(D),\mathscr M\rangle\simeq N_{D/S}(\mathscr M).\]
若 $D,E$ 是\Term{disjoint}{不相交}的\Term{effective}{有效}\Term{relative-divisor}{相对除子}，则 $\mathscr O(E)|D$ \Term{canonical}{典范}地\Term{isomorphism}{同构}于 $\mathscr O$，故有\Term{isomorphism}{同构}
\begin{equation}\sigma_{D,E}:\langle\mathscr O(D),\mathscr O(E)\rangle\simeq N_{D/S}(\mathscr O(E))\simeq N_{D/S}(\mathscr O)\simeq\mathscr O.\tag{1.3.11.1}\end{equation}
\end{SGAresult}
\begin{SGAresult}{命题}{1.3.12}
设 $(D_i)_{i=1,2}$ 和 $E$ 为\Term{effective}{有效}\Term{relative-divisor}{相对除子}，且 $D_i$ 与 $E$ \Term{disjoint}{不相交}（$i=1,2$）。设 $f$ 为满足 $\operatorname{div}(f)=D_1-D_2$ 的\Term{rational-function}{有理函数}，即\Term{isomorphism}{同构} $f:\mathscr O(D_1)\xrightarrow\sim\mathscr O(D_2)$。则以下\Term{isomorphism}{同构}图交换：
\SourcePage{332}
\begin{equation}\begin{tikzcd}\langle\mathscr O(D_1),\mathscr O(E)\rangle\ar[r,"{\sigma_{D_1,E}}"]\ar[d,"{\langle f,\mathscr O(E)\rangle}"']&\mathscr O\ar[d,"{N_{E/S}(f)}"]\\\langle\mathscr O(D_2),\mathscr O(E)\rangle\ar[r,"{\sigma_{D_2,E}}"]&\mathscr O\end{tikzcd}\tag{1.3.12.1}\end{equation}
\begin{equation}\begin{tikzcd}\langle\mathscr O(E),\mathscr O(D_1)\rangle\ar[r,"{\sigma_{E,D_1}}"]\ar[d,"{\langle\mathscr O(E),f\rangle}"']&\mathscr O\ar[d,"{N_{E/S}(f)}"]\\\langle\mathscr O(E),\mathscr O(D_2)\rangle\ar[r,"{\sigma_{E,D_2}}"]&\mathscr O.\end{tikzcd}\tag{1.3.12.2}\end{equation}
\end{SGAresult}
\begin{proof}
公式 (1.3.12.2) 只是表达以下图的交换性：
\[\begin{tikzcd}N_{E/S}(\mathscr O(D_1))\ar[r,"\sim"]\ar[d,"{N_{E/S}(f)}"']&\mathscr O\ar[d,"{N_{E/S}(f)}"]\\N_{E/S}(\mathscr O(D_2))\ar[r,"\sim"]&\mathscr O.\end{tikzcd}\]
对每个 $S$-\Term{scheme}{概形} $T$ 和 $X_T$ 上的\Term{invertible-sheaf}{可逆层} $\mathscr L$，置 $F(\mathscr L)=N_{E_T/T}(\mathscr L)$。此\Term{functor}{函子}满足 1.3.10 的假设，故存在唯一的\Term{invertible-sheaf}{可逆层} $\mathscr M$，带有与双\Term{additive}{加性}相容的\Term{functor}{函子}\Term{isomorphism}{同构}
\[N_{E_T/T}(\mathscr L)\xrightarrow\sim\langle\mathscr L,\mathscr M\rangle.\]
此外，若 $d:T\to X_T$ 为 $X/T$ 的\Term{section}{截面}，则\Term{canonical}{典范}地有 $d^*\mathscr M\xrightarrow\sim N_{E_T/T}(\mathscr O(d))$。取 $d$ 为 $X$ 的通有\Term{section}{截面} $\Delta:X\to X_X$，得到
\[\mathscr M\xrightarrow\sim N_{E_X/X}(\mathscr O(\Delta)).\]
在 $S$ 上\Term{local-scheme}{局部}地，$E$ 是\Term{section}{截面}之和 $E=\sum t_i$（XVII 6.3）；对此种 $E$，有
\[N_{E_X/X}(\mathscr O(\Delta))=\bigotimes t_i^*\mathscr O(\Delta)=\sum\mathscr O(t_i)=\mathscr O(E).\]
由于不依赖 $t_i$ 的次序，此\Term{isomorphism}{同构}定义\Term{canonical-isomorphism}{典范同构}
\begin{equation}N_{E_X/X}(\mathscr O(\Delta))\simeq\mathscr O(E),\tag{1.3.12.3}\end{equation}
故有
\begin{equation}N_{E_T/T}(\mathscr L)\xrightarrow\sim\langle\mathscr L,\mathscr O(E)\rangle.\tag{1.3.12.4}\end{equation}
\SourcePage{333}
\Term{isomorphism}{同构} (1.3.12.3) 与双\Term{additive}{加性}、\Term{base-change}{基变换}相容，且由下列附\Term{additive}{加性}质刻画：对 $X_T/T$ 的每个\Term{section}{截面} $d$，\Term{isomorphism}{同构}
\[N_{E/T}(\mathscr O(d))\xrightarrow\sim\langle\mathscr O(d),\mathscr O(E)\rangle\xrightarrow\sim d^*\mathscr O(E)\]
为通有\Term{isomorphism}{同构} (1.3.12.3) 的\Term{inverse-image}{逆像}。特别地，当 $d$ 与 $E$ \Term{disjoint}{不相交}时，下图交换：
\[\begin{tikzcd}N_{E/T}(\mathscr O(d))\ar[r,"\sim"]\ar[d,equal]&\langle\mathscr O(d),\mathscr O(E)\rangle\ar[r,"\sim"]&d^*\mathscr O(E)\ar[d,equal]\\\mathscr O\ar[rr,equal]&&\mathscr O.\end{tikzcd}\]
由加法推出：当 $D$ 与 $E$ \Term{disjoint}{不相交}时，下图交换：
\begin{equation}\begin{tikzcd}N_{E/T}(\mathscr O(D))\ar[r,"\sim"]\ar[d,equal]&\langle\mathscr O(D),\mathscr O(E)\rangle\ar[r,"\sim"]&N_{D/T}(\mathscr O(E))\ar[d,equal]\\\mathscr O\ar[rr,equal]&&\mathscr O.\end{tikzcd}\tag{1.3.12.5}\end{equation}
所以\Term{arrow}{箭头} $\sigma_{E,D_1}$ 是 1.3.12.4 的逆与一个显然的\Term{arrow}{箭头}的复合：
\[\langle\mathscr O(D),\mathscr O(E)\rangle\xleftarrow\sim N_{E/S}(\mathscr O(D))\longrightarrow\mathscr O.\]
由于 1.3.12.4 为\Term{functorial}{函子性}，(1.3.12.2) 立即得到。
\end{proof}
从 1.3.12 推出（参见 Serre [1], Ch. III, prop. 7, p. 46）：
\begin{SGAresult}{推论}{1.3.13}
设 $(D_i)_{i=1,2}$、$(E_i)_{i=1,2}$ 为\Term{effective}{有效}\Term{relative-divisor}{相对除子}，$D_i$ 与 $E_j$ \Term{disjoint}{不相交}，$f$ 为满足 $\operatorname{div}(f)=D_1-D_2$ 的\Term{rational-function}{有理函数}，$g$ 为满足 $\operatorname{div}(g)=E_1-E_2$ 的\Term{rational-function}{有理函数}。则有
\[N_{E_1/S}(f)N_{E_2/S}(f)^{-1}=N_{D_1/S}(g)N_{D_2/S}(g)^{-1},\]
简写为
\begin{equation}f(\operatorname{div}(g))=g(\operatorname{div}(f)).\tag{1.3.13.1}\end{equation}
\end{SGAresult}
此公式写成 $N_{D_2/S}(g)N_{E_1/S}(f)=N_{E_2/S}(f)N_{D_1/S}(g)$ 的\Term{form}{形式}时，通过 (1.3.12.1)、(1.3.12.2) 表达下图的交换性：
\[\begin{tikzcd}\langle\mathscr O(D_1),\mathscr O(E_1)\rangle\ar[r,"{\langle\mathscr O(D_1),g\rangle}"]\ar[d,"{\langle f,\mathscr O(E_1)\rangle}"']&\langle\mathscr O(D_1),\mathscr O(E_2)\rangle\ar[d,"{\langle f,\mathscr O(E_2)\rangle}"]\\\langle\mathscr O(D_2),\mathscr O(E_1)\rangle\ar[r,"{\langle\mathscr O(D_2),g\rangle}"]&\langle\mathscr O(D_2),\mathscr O(E_2)\rangle.\end{tikzcd}\]
\SourcePage{334}
\begin{SGAresult}{推论}{1.3.14}
当 $G=\mathbf G_m$ 时，次数 (1.3.2.7) 与\Term{invertible-sheaf}{可逆层}的通常次数一致。
\end{SGAresult}
只须对从 $f_0\in\Gamma(S,\mathscr O_S^*)$ \Term{inverse-image}{逆像}而来的 $f$ 比较公式 (1.3.7.3)、(1.3.12.1)。
\textbf{1.3.15.} 对任意充分\Term{ample}{丰沛}的 $\mathscr L,\mathscr M$，在 $S$ 上\Term{local-scheme}{局部}地，存在\Term{disjoint}{不相交}\Term{divisor}{除子} $D,E$ 和\Term{isomorphism}{同构}
\[\alpha:\mathscr L\longrightarrow\mathscr O(D),\qquad\beta:\mathscr M\longrightarrow\mathscr O(E).\]
以 $\tau$ 表示使下图交换的\Term{isomorphism}{同构}：
\[\begin{tikzcd}\langle\mathscr L,\mathscr M\rangle\ar[r,"{\langle\alpha,\beta\rangle}"]\ar[d,"\tau"']&\langle\mathscr O(D),\mathscr O(E)\rangle\ar[r,"{\sigma_{D,E}}"]&\mathscr O\ar[d,equal]\\\langle\mathscr M,\mathscr L\rangle\ar[r,"{\langle\beta,\alpha\rangle}"]&\langle\mathscr O(E),\mathscr O(D)\rangle\ar[r,"{\sigma_{E,D}}"]&\mathscr O.\end{tikzcd}\]
由 1.3.12，此\Term{isomorphism}{同构}不依赖 $D,E,\alpha,\beta$ 的具体选取，故可整体化为\Term{canonical-isomorphism}{典范同构}
\begin{equation}\tau:\langle\mathscr L,\mathscr M\rangle\xrightarrow\sim\langle\mathscr M,\mathscr L\rangle.\tag{1.3.15.1}\end{equation}
此\Term{isomorphism}{同构}满足 $\tau^2=1$；由于 (1.3.11.1) 与双\Term{additive}{加性}相容，它亦如此，并可线性地延拓至任意\Term{invertible-sheaf}{可逆层}。
由 (1.3.12.5)，\Term{morphism}{态射} (1.3.12.4) 等同于复合
\[N_{E_T/T}(\mathscr L)\xleftarrow\sim\langle\mathscr O(E),\mathscr L\rangle\xrightarrow[\sim]{\tau}\langle\mathscr L,\mathscr O(E)\rangle\]
（对与 $E$ \Term{disjoint}{不相交}的 $D$，取 $\mathscr L=\mathscr O(D)$ 验证）。
\begin{SGAresult}{公式汇编}{1.3.16}
对每条\Term{projective}{射影}\Term{smooth}{光滑}曲线 $f:X\to S$，有构造与任意\Term{base-change}{基变换}相容的\Term{functor}{函子} $\langle\ ,\ \rangle$，它将 $X$ 上的两个\Term{invertible-sheaf}{可逆层} $\mathscr L,\mathscr M$ 送到 $S$ 上的\Term{invertible-sheaf}{可逆层} $\langle\mathscr L,\mathscr M\rangle$。此\Term{functor}{函子}带有下列附加结构，并满足下列条件。
\end{SGAresult}
\textbf{1.3.16.1.} 采用 1.3.8 的记号，当 $G=\mathbf G_m$ 时，有 $\langle\mathscr L,\mathscr M\rangle=\langle\mathscr L,\mathscr M]$。
特别地，$\langle\mathscr L,\mathscr M\rangle$ 在 1.3.8.1 指定的意义下对 $\mathscr L,\mathscr M$ 双\Term{additive}{加性}，且此双\Term{additive}{加性}与关于 $\mathscr M$ \Term{functorial}{函子性}的\Term{canonical-isomorphism}{典范同构}
\[\langle\mathscr O(E),\mathscr M\rangle\xrightarrow\sim N_{E/S}(\mathscr M)\]
右侧的双\Term{additive}{加性}相容。
\SourcePage{335}
\textbf{1.3.16.2.} 有\Term{symmetry-isomorphism}{对称同构} $\tau:\langle\mathscr L,\mathscr M\rangle\xrightarrow\sim\langle\mathscr M,\mathscr L\rangle$，它对 $\mathscr L,\mathscr M$ \Term{functorial}{函子性}，与\Term{base-change}{基变换}、双\Term{additive}{加性}\Term{isomorphism}{同构}相容，且满足 $\tau^2=1$。对\Term{disjoint}{不相交}\Term{divisor}{除子} $D,E$，它使下图交换：
\[\begin{tikzcd}[column sep=small]N_{D/S}(\mathscr O(E))\ar[d,equal]&\langle\mathscr O(D),\mathscr O(E)\rangle\ar[l]\ar[r,"\tau"]&\langle\mathscr O(E),\mathscr O(D)\rangle\ar[r]&N_{E/S}(\mathscr O(D))\ar[d,equal]\\\mathscr O\ar[rrr,equal]&&&\mathscr O.\end{tikzcd}\]
这些公式蕴含公式 (1.3.12.1)、(1.3.12.2)、(1.3.13.1)，并且蕴含：
\textbf{1.3.16.3.} 对 $a,b\in\Gamma(S,\mathscr O_S^*)$，以 $\langle a,b\rangle$ 表示由 $\mathscr L,\mathscr M$ 的\Term{automorphism}{自同构} $a,b$ 按\Term{functorial}{函子性}\Term{derive}{导出}的 $\langle\mathscr L,\mathscr M\rangle$ 的\Term{automorphism}{自同构}，并视为 $\mathbf G_m$ 的\Term{section}{截面}，则
\[\langle a,b\rangle=a^{\deg(\mathscr M)}b^{\deg(\mathscr L)}.\]
\textbf{1.3.16.4.} 若 $\mathscr L_0$（相应地 $\mathscr M_0$）为 $S$ 上的\Term{invertible-module}{可逆模}，则有\Term{canonical-isomorphism}{典范同构}
\[\langle f^*\mathscr L_0,\mathscr M\rangle\xrightarrow\sim\mathscr L_0^{\otimes\deg(\mathscr M)}\qquad
(\text{相应地 }\langle\mathscr L,f^*\mathscr M_0\rangle\xrightarrow\sim\mathscr M_0^{\otimes\deg(\mathscr L)}).\]
\textbf{1.3.16.5.} 若 $u:X\to Y$ 为 $S$ 上\Term{projective}{射影}\Term{smooth}{光滑}曲线之间的\Term{finite}{有限}\Term{flat}{平坦}\Term{morphism}{态射}，则对 $X$ 上的\Term{invertible-sheaf}{可逆层} $\mathscr L$ 和 $Y$ 上的\Term{invertible-sheaf}{可逆层} $\mathscr M$，\Term{isomorphism}{同构} 1.2.9.7
\[\langle\mathscr L,u^*\mathscr M\rangle\simeq\langle N_{X/Y}\mathscr L,\mathscr M\rangle,\qquad
\langle u^*\mathscr M,\mathscr L\rangle\simeq\langle\mathscr M,N_{X/Y}\mathscr L\rangle\]
通过对称性 (1.2.16.2) 互相\Term{derive}{导出}（取 $\mathscr L=\mathscr O(E)$，$\mathscr M=\mathscr O(F)$ 且 $E,F$ \Term{disjoint}{不相交}验证）。
\textbf{1.3.16.6.} 对 $X$ 上的\Term{invertible-sheaf}{可逆层} $\mathscr L$，对称\Term{arrow}{箭头} $\tau:\langle\mathscr L,\mathscr L\rangle\to\langle\mathscr L,\mathscr L\rangle$ 是乘以 $(-1)^{\deg(\mathscr L)}$。
\textbf{1.3.17.} 剩下证明 1.3.16.6。我们给出两个证明，第一个使用行列式理论（SGA 6 XI）。
\textbf{第一个证明。} 设 $u:T\to S$ 为\Term{finite}{有限}\Term{locally-free}{局部自由}\Term{morphism}{态射}。若 $\mathscr L$ 为 $T$ 上的\Term{invertible-sheaf}{可逆层}，则有\Term{isomorphism}{同构}
\begin{equation}N_{T/S}(\mathscr L)\simeq\det(u_*\mathscr L)\det(u_*\mathscr O_T)^{-1},\tag{1.3.17.1}\end{equation}
由以下三个条件唯一确定：
\SourcePage{336}
a) 其构造与\Term{base-change}{基变换}相容；b) 当 $\mathscr L=\mathscr O_T$ 时，它是 $\mathscr O_S$ 的恒等\Term{morphism}{态射}；c) 它对 $\mathscr L$ \Term{functorial}{函子性}。
设 $D,E$ 为 $X$ 上的两个\Term{effective}{有效}\Term{relative-divisor}{相对除子}，$i$ 为 $D$ 到 $X$ 的包含，$u$ 为 $D$ 到 $S$ 的\Term{projection}{投影}。短\Term{exact-sequence}{正合序列}
\[0\longrightarrow\mathscr O(-E)\longrightarrow\mathscr O\longrightarrow\mathscr O_E\longrightarrow0\]
允许将\Term{isomorphism}{同构}
\[\langle\mathscr O(D),\mathscr O(-E)\rangle\simeq N_{D/S}(\mathscr O(-E))\simeq\det(u_*i^*\mathscr O(-E))\det(u_*i^*\mathscr O)^{-1}\]
解释为\Term{isomorphism}{同构}
\[\langle\mathscr O(D),\mathscr O(-E)\rangle=\det\mathrm Ru_*(Li^*\mathscr O_E)^{-1}\simeq\det\mathrm Rf_*(\mathscr O_D\mathbin{\mathop\otimes^{\mathrm L}}\mathscr O_E)^{-1},\]
或
\begin{equation}\langle\mathscr O(D),\mathscr O(E)\rangle\simeq\det\mathrm Rf_*(\mathscr O_D\mathbin{\mathop\otimes^{\mathrm L}}\mathscr O_E).\tag{1.3.17.2}\end{equation}
当 $D,E$ \Term{disjoint}{不相交}时，有 $\mathscr O_D\mathbin{\mathop\otimes^{\mathrm L}}\mathscr O_E=0$，由此得到的 $\langle\mathscr O(D),\mathscr O(E)\rangle$ 与 $\mathscr O_S$ 的\Term{isomorphism}{同构}，就是先前考虑的\Term{isomorphism}{同构}。因此当 $D,E$ \Term{disjoint}{不相交}时，对称性 $\tau$ 与\Term{symmetry-isomorphism}{对称同构}
\[\mathscr O_D\mathbin{\mathop\otimes^{\mathrm L}}\mathscr O_E\xrightarrow\sim\mathscr O_E\mathbin{\mathop\otimes^{\mathrm L}}\mathscr O_D\]
通过 (1.3.17.2) 相容。经\Term{specialization}{特殊化}，对任意 $D,E$ 仍如此。
当 $D=E$ 由\Term{ideal}{理想} $I$ 定义时，有
\[\begin{array}{lll}
H^0(\mathscr O_D\mathbin{\mathop\otimes^{\mathrm L}}\mathscr O_D)=\mathscr O_D,&\text{对称性}=+1,\\
H^{-1}(\mathscr O_D\mathbin{\mathop\otimes^{\mathrm L}}\mathscr O_D)=I/I^2,&\text{对称性}=-1,\\
H^i(\mathscr O_D\mathbin{\mathop\otimes^{\mathrm L}}\mathscr O_D)=0,&i\ne0,-1.
\end{array}\]
于是\Term{arrow}{箭头} $\tau:\langle\mathscr O(D),\mathscr O(D)\rangle\xrightarrow\sim\langle\mathscr O(D),\mathscr O(D)\rangle$ 是 $(-1)^{\dim I/I^2}=(-1)^{\deg(\mathscr O(D))}$，1.3.16.6 由此得到。
\textbf{第二个证明。} 归结于“通有”情形，可假设 $S$ \Term{reduced}{约化}。此情形又归结于 $S$ 为\Term{algebraically-closed-field}{代数闭域} $k$ 的谱。此时设 $X_1,X_2,Y_1,Y_2$ 为 $X$ 上的四个\Term{divisor}{除子}，$f,g$ 为两个\Term{rational-function}{有理函数}，并假设
\[\operatorname{div}(f)=X_2-X_1,\qquad\operatorname{div}(g)=Y_2-Y_1,\qquad(X_1\cup Y_2)\cap(X_2\cup Y_1)=\varnothing.\]
于是映射 $f,g$ 定义
\[\langle f,g\rangle:\langle\mathscr O(X_1),\mathscr O(Y_1)\rangle\longrightarrow\langle\mathscr O(X_2),\mathscr O(Y_2)\rangle.\]
\SourcePage{337}
两侧均\Term{canonical}{典范}地\Term{isomorphism}{同构}于 $k$（1.3.11.1），故 $\langle f,g\rangle$ 可视为 $k^*$ 中的元素。
\begin{SGAresult}{引理}{1.3.18}
有
\[\langle f,g\rangle=\prod_{x\in X_1\cup Y_2}(-1)^{v_x(f)v_x(g)}(g^{v_x(f)}/f^{v_x(g)})(x).\]
\end{SGAresult}
先证明 (1.3.18) 蕴含 (1.3.16.6)。可假设 $\mathscr L\simeq\mathscr O(X)\simeq\mathscr O(Y)$，其中 $X,Y$ \Term{effective}{有效}，且 $X\cap Y=\varnothing$。设 $f$ 为满足 $\operatorname{div}(f)=Y-X$ 的\Term{rational-function}{有理函数}。\Term{commutative-diagram}{交换图}
\[\begin{tikzcd}\langle\mathscr L,\mathscr L\rangle\ar[r,"\tau"]\ar[d,equal]&\langle\mathscr L,\mathscr L\rangle\ar[d,equal]\\\langle\mathscr O(X),\mathscr O(Y)\rangle\ar[r,"{\langle f,f^{-1}\rangle}"]&\langle\mathscr O(Y),\mathscr O(X)\rangle\end{tikzcd}\]
表明
\[\tau=\langle f,f^{-1}\rangle=\prod_{x\in X}(-1)^{v_x(f)^2}=(-1)^{\deg(\mathscr L)}.\]
下面证明 1.3.18。可归结于 $X_i,Y_i$ 很\Term{ample}{丰沛}的情形。于是存在\Term{resolution}{分解} $f=f_1f_2$、$g=g_1g_2$，使
\[\begin{aligned}\operatorname{div}(f_1)&=X'-X_1,&\operatorname{div}(g_1)&=Y'-Y_1,\\\operatorname{div}(f_2)&=X_2-X',&\operatorname{div}(g_2)&=Y_2-Y',\end{aligned}\]
其中\Term{divisor}{除子} $X',Y'$ 和 $X_1\cup X_2\cup Y_1\cup Y_2$ 互\Term{disjoint}{不相交}。
以 $\langle\ ,\ \rangle^*$ 表示 (1.3.18) 的右侧，并证明
\[\langle f,g\rangle^*=\langle f_1,g_1\rangle^*\langle f_2,g_2\rangle^*.\]
由 Serre [1]，有
\begin{equation}\prod_x(-1)^{v_x(f_1)v_x(g_2)}(f_1^{v_x(g_2)}/g_2^{v_x(f_1)})(x)=1.\tag{1.3.18.1}\end{equation}
由此得到
\[\begin{aligned}
\langle f_1,g_1\rangle^*\langle f_2,g_2\rangle^*
={}&g_1(X_1)f_1(Y')g_2(X')f_2(Y_2)\\
&\cdot\prod_{X_1\cup X'\cup Y'\cup Y_2}(-1)^{v_x(f_1)v_x(g_2)}(f_1^{v_x(g_2)}f_2^{-v_x(f_1)})(x)\\
={}&\prod_{x\in X'}g_2(x)^{v_x(f_1)}g_2(x)^{-v_x(f_1)}
\prod_{x\in Y'}f_1(x)^{v_x(g_1)}f_1(x)^{v_x(g_2)}\\
&\cdot\prod_{X_1\cup Y_2}(-1)^{v_x(f)v_x(g)}((f_1f_2)^{v_x(g_2)}(g_1g_2)^{-v_x(f_1)})(x)\\
={}&\langle f_1f_2,g_1g_2\rangle^*.
\end{aligned}\]
由 $\langle\ ,\ \rangle$ 的\Term{functorial}{函子性}，有 $\langle f,g\rangle=\langle f_1,g_1\rangle\langle f_2,g_2\rangle$，而容易验证 $\langle f_i,g_i\rangle=\langle f_i,g_i\rangle^*$。引理由此得证。
